\documentclass[11pt]{article}

\usepackage[T1]{fontenc}
\usepackage[utf8]{inputenc}
\usepackage{lmodern}
\usepackage[margin=1.06in]{geometry}
\usepackage{amsmath,amssymb,mathtools,amsthm}
\usepackage{aliascnt}
\usepackage{mathrsfs}
\usepackage{microtype}
\usepackage{enumitem}
\usepackage{xcolor}
\usepackage{url}
\usepackage{hyperref}

\hypersetup{
  colorlinks=true,
  linkcolor=blue!45!black,
  citecolor=green!35!black,
  urlcolor=blue!55!black,
  pdftitle={The Two-Variable Factorial Conjecture},
  pdfauthor={Christopher D. Long},
  pdfsubject={Factorial Conjecture, Mittag-Leffler constellations, Fourier-Laplace rigidity},
  pdfkeywords={Factorial Conjecture, polynomial moments, Mittag-Leffler function, constellation monodromy, E-function, Fourier-Laplace transform}
}

\numberwithin{equation}{section}
\setlength{\emergencystretch}{3em}

\newtheorem{theorem}{Theorem}[section]

\newaliascnt{proposition}{theorem}
\newtheorem{proposition}[proposition]{Proposition}
\aliascntresetthe{proposition}

\newaliascnt{lemma}{theorem}
\newtheorem{lemma}[lemma]{Lemma}
\aliascntresetthe{lemma}

\newaliascnt{corollary}{theorem}
\newtheorem{corollary}[corollary]{Corollary}
\aliascntresetthe{corollary}

\theoremstyle{definition}
\newaliascnt{definition}{theorem}
\newtheorem{definition}[definition]{Definition}
\aliascntresetthe{definition}

\theoremstyle{remark}
\newaliascnt{remark}{theorem}
\newtheorem{remark}[remark]{Remark}
\aliascntresetthe{remark}

% Load cleveref after the theorem environments so shared counters retain
% their proposition/lemma/corollary names in cross-references.
\usepackage[capitalize,noabbrev]{cleveref}

\newcommand{\C}{\mathbb C}
\newcommand{\Q}{\mathbb Q}
\newcommand{\Qbar}{\overline{\mathbb Q}}
\newcommand{\N}{\mathbb N}
\newcommand{\Z}{\mathbb Z}
\newcommand{\Aone}{\mathbb A^1}
\newcommand{\Gm}{\mathbb G_m}
\newcommand{\Lfac}{\mathcal L}
\newcommand{\FC}{\mathrm{FC}}
\newcommand{\Span}{\operatorname{span}}
\newcommand{\Tr}{\operatorname{Tr}}
\newcommand{\Spec}{\operatorname{Spec}}
\newcommand{\End}{\operatorname{End}}
\newcommand{\ord}{\operatorname{ord}}
\newcommand{\diag}{\operatorname{diag}}
\newcommand{\Bcal}{\mathcal B}
\newcommand{\Gcal}{\mathcal G}
\newcommand{\Jcal}{\mathcal J}
\newcommand{\Kcal}{\mathcal K}
\newcommand{\Mcal}{\mathcal M}
\newcommand{\Ecal}{\mathcal E}
\newcommand{\Rcal}{\mathcal R}
\newcommand{\Ncal}{\mathcal N}
\newcommand{\Lcal}{\mathscr L}
\newcommand{\Vcal}{\mathscr V}
\newcommand{\Dmod}{\mathcal D}
\newcommand{\Ocal}{\mathcal O}
\newcommand{\eps}{\varepsilon}
\newcommand{\dd}{\,d}

\title{The Two-Variable Factorial Conjecture\\
via Mittag--Leffler Constellations and Fourier--Laplace Rigidity}
\author{Christopher D. Long}
\date{Revised research draft, August 1, 2026}

\begin{document}

\maketitle

\begin{center}
\small
Headlamp Software, Lexington, Kentucky, USA\\
\texttt{galizur@gmail.com}
\end{center}

\begin{abstract}
Let $\Lfac(x^a y^b)=a!b!$ on $\C[x,y]$.  We prove the two-variable
Factorial Conjecture: if $\Lfac(f^m)=0$ for every $m\ge1$, then $f=0$.
The proof has three parts.  First, algebraic specialization reduces a
hypothetical counterexample to coefficients in $\Qbar$.  Second, an exact
Mittag--Leffler generating period rules out every nonflat projective leading
form.  Its parameterized radial reversion has a computable logarithmic
symbol; the positive-interval constellation of Pakovich and Muzychuk forces
a primitive Fourier relation, and inverse-branch asymptotics at infinity
contradict that relation.  In the remaining flat-top case
$f_D=(x+y)^D$, the leading singular coefficient of the normalized generating
function is
\[
  \int_0^1\exp\!\left(\frac{f_{D-1}(t,1-t)}D\right)\,dt,
\]
so vanishing of all factorial moments forces this integral to vanish.
Finally, we prove that for every nonconstant $q\in\Qbar[x]$ and distinct
$A,B\in\Qbar$, the integral $\int_A^B e^{q(x)}\,dx$ is transcendental.
The arithmetic proof uses an inhomogeneous $E$-function moment system,
semisimplicity of its Fourier--Laplace connection, horizontal-endomorphism
rigidity in the Jacobian algebra of $q$, Beukers' refinement of the
Siegel--Shidlovskii theorem, and a Cauchy-transform moment argument.  Applied
to $q(t)=f_{D-1}(t,1-t)/D$, this gives the final contradiction.
\end{abstract}

\tableofcontents

\section{Introduction}\label{sec:introduction}

For $g\in\C[x,y]$, define the factorial functional by
\begin{equation}\label{eq:def-L}
  \Lfac(x^a y^b)=a!b!,\qquad a,b\ge0.
\end{equation}
The Factorial Conjecture was introduced by van den Essen, Wright, and Zhao
in their study of the Image Conjecture \cite[\S4]{vdEWZ}.  Its two-variable form
is the following.

\begin{theorem}[Two-variable Factorial Conjecture]\label{thm:FC2}
For every $f\in\C[x,y]$,
\[
  \Lfac(f^m)=0\quad\text{for all }m\ge1
  \qquad\Longrightarrow\qquad f=0.
\]
\end{theorem}

The homogeneous two-variable case was proved by Liu and Sun
\cite{LiuSun}.  The inhomogeneous problem is not formally reduced to that
case because radial integration couples the homogeneous layers through
Gamma factors.

Write a nonzero polynomial of exact degree $D$ as
\[
  f=f_D+f_{D-1}+\cdots+f_0,
\]
where every $f_j$ is homogeneous of degree $j$, and put
\begin{equation}\label{eq:intro-projective-layers}
  A_j(t)=f_{D-j}(t,1-t),\qquad 0\le j\le D.
\end{equation}
The proof of \cref{thm:FC2} follows the implication chain
\[
\begin{aligned}
\text{counterexample over }\C
&\Longrightarrow \text{counterexample over }\Qbar,\\
&\Longrightarrow A_0(t)\text{ is constant},\\
&\Longrightarrow f_D=(x+y)^D\text{ after scaling},\\
&\Longrightarrow \int_0^1 e^{A_1(t)/D}\,dt=0,\\
&\Longrightarrow \text{contradiction}.
\end{aligned}
\]
The second implication is the nonflat-top theorem proved in
\cref{sec:radial,sec:mittag,sec:symbol,sec:constellation,sec:infinity}.
The fourth is the flat first-symbol calculation in \cref{sec:flat-symbol}.
The final contradiction is supplied by the algebraic exponential-integral
theorem of \cref{thm:exp-integral}, proved in
\cref{sec:E-functions,sec:exp-moment,sec:exp-semisimple,sec:exp-rigidity,sec:exp-relations,sec:exp-annihilation,sec:exp-phase-rigidity}.

\section{Algebraic specialization}\label{sec:specialization}

\begin{proposition}[Algebraic specialization]\label{prop:algebraic-specialization}
If a nonzero counterexample to \cref{thm:FC2} exists in $\C[x,y]$, then
there is one in $\Qbar[x,y]$ with the same monomial support.
\end{proposition}

\begin{proof}
Let
\[
  f=\sum_{\alpha\in S}c_\alpha x^{\alpha_1}y^{\alpha_2}
\]
be a nonzero counterexample with finite support $S$, and choose
$\alpha_0\in S$.  Replacing $f$ by $c_{\alpha_0}^{-1}f$ preserves every
moment vanishing, since
\[
  \Lfac\!\left((c_{\alpha_0}^{-1}f)^n\right)
  =c_{\alpha_0}^{-n}\Lfac(f^n)=0
  \qquad(n\ge1),
\]
and it preserves the support.  We may therefore impose the normalization
$c_{\alpha_0}=1$.

For every $n\ge1$, the quantity $\Lfac(f^n)$ is a polynomial with integer
coefficients in the remaining coefficients $c_\alpha$: the multinomial
coefficients are integers, and applying $\Lfac$ to each resulting monomial
multiplies its coefficient by two factorials.  Put
\[
  R=\Qbar[c_\alpha,u_\alpha:
       \alpha\in S\setminus\{\alpha_0\}],
\]
and let $I\subseteq R$ be generated by all moment polynomials
$\Lfac(f^n)$, $n\ge1$, together with the Rabinowitsch equations
\[
  u_\alpha c_\alpha-1=0
  \qquad(\alpha\in S\setminus\{\alpha_0\}).
\]
Although the displayed family of moment generators is infinite, the ideal
is proper: evaluation at the original normalized complex coefficient vector,
with $u_\alpha=c_\alpha^{-1}$, annihilates every generator.  Hence $I$ is
contained in a maximal ideal $\mathfrak m$ of $R$.  By the weak
Nullstellensatz, $\mathfrak m$ is the vanishing ideal of a point with
coordinates in $\Qbar$.  At that point the Rabinowitsch equations force every
coefficient indexed by $S$ to remain nonzero, while all the factorial-moment
polynomials vanish.  The resulting polynomial in $\Qbar[x,y]$ therefore has
support exactly $S$ and satisfies every factorial-moment equation.
\end{proof}

\section{Radial--projective coordinates}\label{sec:radial}

\begin{lemma}[Gamma integral]\label{lem:factorial-integral}
For every $g\in\C[x,y]$,
\begin{equation}\label{eq:factorial-integral}
  \Lfac(g)
  =\int_0^\infty\int_0^\infty
    g(x,y)e^{-x-y}\,dx\,dy.
\end{equation}
\end{lemma}

\begin{proof}
The integral is absolutely convergent for every polynomial, since a
polynomial has at most polynomial growth and the weight is $e^{-x-y}$.
It is therefore enough by linearity to test the monomial $x^a y^b$.
Fubini's theorem and the Gamma integral give
\[
  \int_0^\infty\int_0^\infty
  x^a y^b e^{-x-y}\,dx\,dy
  =\Gamma(a+1)\Gamma(b+1)=a!b!.
\]
\end{proof}

Let $f$ have exact degree $D$, with homogeneous decomposition and
projective layers $A_j$ as in \eqref{eq:intro-projective-layers}.

\begin{lemma}[Projective degree and the factor $x+y$]\label{lem:projective-degree}
Let $f_D$ be homogeneous of degree $D$, and write
\[
  f_D=(x+y)^m g_D,
  \qquad x+y\nmid g_D.
\]
Then
\[
  \deg f_D(t,1-t)=D-m.
\]
In particular, $f_D(t,1-t)$ is constant if and only if $f_D$ is a scalar multiple of $(x+y)^D$.
\end{lemma}

\begin{proof}
Since $t+(1-t)=1$,
\[
  f_D(t,1-t)=g_D(t,1-t).
\]
The polynomial on the right has degree at most $D-m$.  Its coefficient of
$t^{D-m}$ is $g_D(1,-1)$.  For a homogeneous polynomial $g_D$, the equality
$g_D(1,-1)=0$ is equivalent to the identity $g_D(x,-x)=0$, and hence to the
divisibility $x+y\mid g_D$.  The defining assumption therefore gives
$g_D(1,-1)\ne0$, so the degree is exactly $D-m$.  The final assertion follows
by taking $D-m=0$.
\end{proof}

For $r>0$, homogeneity gives
\begin{equation}\label{eq:radial-polynomial}
  f(rt,r(1-t))
  =r^D\left(A_0(t)+\frac{A_1(t)}r+\cdots+\frac{A_D(t)}{r^D}\right).
\end{equation}

\begin{proposition}[Exact radial--projective formula]\label{prop:radial-projective}
For every $n\ge0$,
\begin{equation}\label{eq:normalized-radial}
\begin{split}
  I_n(f)&:=\frac{\Lfac(f^n)}{\Gamma(Dn+2)}\\
  &=\frac1{\Gamma(Dn+2)}
    \int_0^1\!\int_0^\infty
    e^{-r}r\,f(rt,r(1-t))^n\,dr\,dt\\
  &=\frac1{\Gamma(Dn+2)}
    \int_0^1\!\int_0^\infty
    e^{-r}r^{Dn+1}
    \left(A_0(t)+\frac{A_1(t)}r+\cdots+\frac{A_D(t)}{r^D}\right)^n
    \,dr\,dt.
\end{split}
\end{equation}
In the last integral the apparent singularity at $r=0$ is removable: the
full integrand is the preceding nonsingular expression
$e^{-r}r\,f(rt,r(1-t))^n$.
\end{proposition}

\begin{proof}
Use the change of variables
\[
  x=rt,\qquad y=r(1-t),
  \qquad r\ge0,\ 0\le t\le1.
\]
The absolute value of its Jacobian determinant is $r$, and $x+y=r$.
Substitution into \eqref{eq:factorial-integral} gives the first integral in
\eqref{eq:normalized-radial}.  For $r>0$, insert
\eqref{eq:radial-polynomial}; the resulting equality extends across $r=0$
by the removable-singularity observation in the statement.
\end{proof}


\section{Mittag--Leffler generating periods}\label{sec:mittag}

Throughout \cref{sec:mittag,sec:symbol,sec:constellation,sec:infinity}, fix $D\ge1$ and polynomials
\[
  A_0,A_1,\ldots,A_D\in\C[t].
\]
Put
\begin{equation}\label{eq:P-general}
  P(t,r)=\sum_{j=0}^D A_j(t)r^{D-j}.
\end{equation}
For the projective layers of a degree-$D$ polynomial, this is exactly $f(rt,r(1-t))$.

For $\beta\in\N$, define the Mittag--Leffler function
\begin{equation}\label{eq:mittag-general}
  E_{D,\beta}(z)=\sum_{n\ge0}\frac{z^n}{\Gamma(Dn+\beta)}.
\end{equation}

\begin{lemma}[Uniform Mittag--Leffler growth and initial holomorphy]
\label{lem:mittag-growth}
For $\beta\in\{2,3\}$ there are constants $C,N>0$, depending only on
$D$ and $\beta$, such that
\begin{equation}\label{eq:mittag-growth}
  |E_{D,\beta}(z)|
  \le C(1+|z|)^N\exp\!\bigl(C|z|^{1/D}\bigr)
  \qquad(z\in\C).
\end{equation}
Let $K\subset\C$ be compact.  Then there is $C_K>0$ such that
\begin{equation}\label{eq:P-growth}
  |P(t,r)|^{1/D}\le C_K(1+r)
  \qquad(t\in K,\ r\ge0).
\end{equation}
Consequently, for $|s|$ sufficiently small, the integrals in
\eqref{eq:Gcal-Mittag-general} and \eqref{eq:Jcal-Mittag} converge
absolutely and locally uniformly in $t$ on $K$ and in $s$.  They define
holomorphic functions of $s$, and of $t$ whenever $t$ is retained as a
complex parameter; termwise integration of the Mittag--Leffler series is
valid.
\end{lemma}

\begin{proof}
For $|z|\le1$, the entire functions $E_{D,\beta}$ are bounded.  If
$|z|\ge1$, choose $X$ with $X^D=z$.  The elementary root-of-unity filter, recorded precisely in
Lemma~\ref{lem:root-filter-general}, writes
$E_{D,\beta}(X^D)$ as a finite sum of terms of the form
$X^{1-\beta}e^{\omega^kX}$ together with a finite Laurent correction.
Since $|X|=|z|^{1/D}$, this gives \eqref{eq:mittag-growth} after enlarging
$C$ and $N$.

Because $P(t,r)$ has degree at most $D$ in $r$ and its coefficients are
bounded on $K$, one has $|P(t,r)|\le C_K^D(1+r)^D$, which is
\eqref{eq:P-growth}.  Hence the absolute value of either integrand is
bounded by
\[
  C'e^{-r}r(1+r)^{N'}
  \exp\!\bigl(C'C_K|s|^{1/D}(1+r)\bigr).
\]
For $|s|$ small enough this is dominated by
$C''e^{-r/2}(1+r)^{N'+1}$, uniformly for $t\in K$.  Dominated convergence
and the Weierstrass theorem now give all asserted holomorphy and
termwise-integration statements.
\end{proof}

\begin{definition}[Normalized generating periods]\label{def:generating-periods}
Let $I_n=I_n(f)$ denote the normalized radial moment in \eqref{eq:normalized-radial}.  Define
\begin{align}
  \Gcal(s)&=\sum_{n\ge0}I_n s^n,\label{eq:Gcal-general}\\
  \Jcal(s)&=\sum_{n\ge0}\frac{I_n}{Dn+2}s^n,\label{eq:Jcal-series-new}\\
  \Kcal(w)&=\frac1w\Jcal(w^{-1}).\label{eq:Kcal-new}
\end{align}
\end{definition}

\begin{proposition}[Exact Mittag--Leffler representations]\label{prop:Mittag-representations}
For $|s|$ sufficiently small,
\begin{align}
  \Gcal(s)
  &=\int_0^1\int_0^\infty
    e^{-r}r\,E_{D,2}\bigl(sP(t,r)\bigr)\,dr\,dt,\label{eq:Gcal-Mittag-general}\\
  \Jcal(s)
  &=\int_0^1\int_0^\infty
    e^{-r}r\,E_{D,3}\bigl(sP(t,r)\bigr)\,dr\,dt.\label{eq:Jcal-Mittag}
\end{align}
Consequently, for $|w|$ sufficiently large,
\begin{equation}\label{eq:Bcal-Mittag}
  \Kcal(w)=\int_0^1\Bcal(w,t)\,dt,
  \qquad
  \Bcal(w,t)=\frac1w\int_0^\infty
  e^{-r}r\,E_{D,3}\!\left(\frac{P(t,r)}w\right)dr.
\end{equation}
\end{proposition}

\begin{proof}
Expand the Mittag--Leffler functions and use
\eqref{eq:normalized-radial}.  The locally uniform integrable majorant
from Lemma~\ref{lem:mittag-growth} justifies termwise integration and proves
both identities.  Substituting $s=w^{-1}$ and multiplying by $w^{-1}$
gives \eqref{eq:Bcal-Mittag}.
\end{proof}


\begin{lemma}[Root-of-unity filter with low-degree corrections]\label{lem:root-filter-general}
Let $m\ge0$, let $\omega=e^{2\pi i/D}$, and define the finite Laurent polynomial
\[
  C_{D,m}(X)
  =\sum_{\substack{0\le \ell<m\\ \ell\equiv m\pmod D}}
   \frac{X^{\ell-m}}{\ell!}.
\]
Then
\begin{equation}\label{eq:root-filter-general}
  E_{D,m+1}(X^D)
  =\frac1{DX^m}\sum_{k=0}^{D-1}
   \omega^{-mk}e^{\omega^kX}-C_{D,m}(X).
\end{equation}
In particular, $C_{D,1}=0$ unless $D=1$, while $C_{D,2}=0$ for $D\ge3$, $C_{2,2}(X)=X^{-2}$, and
\[
  C_{1,2}(X)=X^{-2}+X^{-1}.
\]
\end{lemma}

\begin{proof}
Expand the exponentials.  The root-of-unity sum retains precisely the indices $\ell\equiv m\pmod D$:
\[
  \frac1{DX^m}\sum_{k=0}^{D-1}\omega^{-mk}e^{\omega^kX}
  =\sum_{\substack{\ell\ge0\\\ell\equiv m\pmod D}}
   \frac{X^{\ell-m}}{\ell!}.
\]
The terms with $\ell\ge m$ form $E_{D,m+1}(X^D)$; the finitely many terms with $\ell<m$ form $C_{D,m}$.
\end{proof}

\section{Parameterized radial reversion and the local logarithmic symbol}\label{sec:symbol}

The local singularity of \eqref{eq:Bcal-Mittag} is controlled entirely by the behavior of $P(t,r)$ at $r=\infty$.

\begin{lemma}[Uniform radial reversion near a nonzero leading value]\label{lem:uniform-nonflat-reversion}
Let $t_0\in\C$ satisfy $u_0=A_0(t_0)\neq0$.  After shrinking a neighborhood $U$ of $t_0$, there is a common exterior region in which
\begin{equation}\label{eq:Phi-nonflat}
  \Phi_t(r)=\left(\frac{P(t,r)}{A_0(t)}\right)^{1/D}
  =r+\frac{A_1(t)}{DA_0(t)}+O(r^{-1})
\end{equation}
has a holomorphic inverse
\begin{equation}\label{eq:chi-nonflat}
  r=\chi(t,Y)
  =Y-\frac{A_1(t)}{DA_0(t)}+O(Y^{-1}).
\end{equation}
All expansions are locally uniform in $t\in U$, and their coefficients are holomorphic in $t$.
\end{lemma}

\begin{proof}
The quotient $P(t,r)/(A_0(t)r^D)$ is $1+O(r^{-1})$ uniformly on compact subsets of $U$.  Choose the $D$-th root asymptotic to $1$, giving \eqref{eq:Phi-nonflat}.  In the coordinate $r^{-1}$, the derivative of $\Phi_t$ at infinity is $1$, uniformly in $t$.  The parameter-dependent holomorphic inverse-function theorem gives a common inverse and the expansion \eqref{eq:chi-nonflat}.
\end{proof}

\begin{lemma}[Uniform exterior-ray deformation]\label{lem:nonflat-common-ray}
Under the hypotheses of Lemma~\ref{lem:uniform-nonflat-reversion}, after
shrinking $U$ and increasing $R$, there are a common exterior sector
$\mathfrak S$, a point $Y_0\in\mathfrak S$, and a ray
\[
  \ell=Y_0+e^{i\theta}[0,\infty)\subset\mathfrak S
\]
with the following property.  For every $t\in U$, the image curve
\[
  \Gamma_t=\Phi_t([R,\infty))
\]
lies in $\mathfrak S$, remains in the common domain of the inverse
$\chi(t,Y)$, and is asymptotic to $\ell$.

Let $F(t,Y)$ be jointly holomorphic on $U\times\mathfrak S$ and satisfy
$F(t,Y)=O(Y^{-1})$ there, locally uniformly in $t$.  For $\eta$ in a
sufficiently small sector such that
\[
  \Re(\eta e^{i\theta})\ge c|\eta|
\]
for some $c>0$, the difference
\begin{equation}\label{eq:common-ray-difference}
  \int_{\Gamma_t}e^{-\eta Y}F(t,Y)\,dY
  -\int_{\ell}e^{-\eta Y}F(t,Y)\,dY
\end{equation}
extends jointly holomorphically to $(t,\eta)\in U\times\{|\eta|<\eps\}$.
\end{lemma}

\begin{proof}
The expansions in Lemma~\ref{lem:uniform-nonflat-reversion} are uniform in
$t$.  Thus, after increasing $R$, the curves $\Gamma_t$ lie in one narrow
exterior sector and admit compatible parameterizations over the ray $\ell$
for which the transverse displacement is $O(1)$ and the tangent difference
is $O(|Y|^{-2})$, locally uniformly in $t$.  Shrinking $U$ if necessary,
join the initial point of each $\Gamma_t$ to $Y_0$ by a path in a fixed
compact annular subset of $\mathfrak S$, depending holomorphically on $t$.

Truncate $\Gamma_t$ and $\ell$ at the same large ray parameter $M$ and join
the corresponding outer endpoints inside $\mathfrak S$.  Their distance,
and hence the length of this outer connector, is $O(1)$ rather than $O(M)$.
On the connector one has $F(t,Y)=O(M^{-1})$.  Moreover, after narrowing the
$\eta$-sector and the exterior sector, $e^{-\eta Y}$ is uniformly bounded
there on every closed admissible subsector together with its vertex
$\eta=0$.  The outer-connector integral is therefore $O(M^{-1})$ uniformly
on compact subsets of $U$ and of such a closed subsector, and tends to zero.

Cauchy's theorem now identifies \eqref{eq:common-ray-difference}, throughout
the admissible sector, with the finite initial-connector integral.  That
integral is entire in $\eta$ and holomorphic in $t$, and therefore furnishes
the asserted joint extension to a full disk about $\eta=0$.  The same
argument applies after termwise differentiation on smaller compact sets.
\end{proof}

\begin{lemma}[A parameterized Laplace logarithm]\label{lem:parameter-Laplace-log}
Let $q_0(t)$ be holomorphic on a neighborhood $U$, and let $H(t,z)$ be holomorphic for $(t,z)\in U\times\{|z|<\rho\}$.  On a fixed sector
\[
  0<|\eta|<\eps,
  \qquad
  |\arg\eta-\theta|<\frac\pi2-\delta,
\]
consider
\[
  L(t,\eta)=\int_{Y_0}^{\infty e^{i\theta}}
  e^{-\eta Y}\left(\frac{q_0(t)}Y+\frac{H(t,Y^{-1})}{Y^2}\right)dY.
\]
Then
\begin{equation}\label{eq:Laplace-log-structure}
  L(t,\eta)=-q_0(t)\log\eta+H_0(t,\eta)
  +\eta\log\eta\,H_1(t,\eta),
\end{equation}
where $H_0,H_1$ are jointly holomorphic in $t$ and $\eta$ near $U\times\{0\}$.
\end{lemma}

\begin{proof}
The $Y^{-1}$ term gives the exponential integral and hence $-\log\eta$ modulo a holomorphic function.  Expand
\[
  H(t,z)=\sum_{m\ge0}h_m(t)z^m
\]
normally on a smaller bidisk.  The Laplace transform of $Y^{-m-2}$ is, modulo a holomorphic function, a nonzero scalar multiple of $\eta^{m+1}\log\eta$.  Normal convergence, uniform on compact subsets of $U$, permits termwise integration and gives \eqref{eq:Laplace-log-structure}.  Equivalently, one may subtract finitely many terms and integrate the analytic remainder repeatedly by parts; the Cauchy estimates for $H$ give the required uniform bounds.
\end{proof}

\begin{theorem}[Parameterized local logarithmic symbol]\label{thm:boundary-symbol}
Let $t_0\in\C$ satisfy $u_0=A_0(t_0)\neq0$.  The germ \eqref{eq:Bcal-Mittag}, initially defined for large $w$, admits multivalued analytic continuation to a neighborhood of $(w,t)=(u_0,t_0)$ outside the divisor $w=A_0(t)$.  On a local slit neighborhood,
\begin{align}\label{eq:boundary-symbol}
  \Bcal(w,t)
  &=-\frac1{DA_0(t)}
  \exp\!\left(\frac{A_1(t)}{DA_0(t)}\right)
  \log\bigl(w-A_0(t)\bigr)\notag\\
  &\quad+H_0(w,t)
  +\bigl(w-A_0(t)\bigr)
   \log\bigl(w-A_0(t)\bigr)H_1(w,t),
\end{align}
where $H_0,H_1$ are jointly holomorphic.  Thus the lowest logarithmic symbol depends only on $A_0$ and $A_1$.
\end{theorem}

\begin{proof}
Shrink the neighborhood $U$ from \cref{lem:uniform-nonflat-reversion} so that $A_0$ never vanishes.  Split the radial integral in \eqref{eq:Bcal-Mittag} at a sufficiently large $R$, increasing $R$ so that $\Phi_t(r)$ is uniformly bounded away from zero for $t\in U$ and $r\ge R$.  The compact part $0\le r\le R$ is holomorphic for $w$ near $u_0\neq0$ because $E_{D,3}$ is entire.

Choose the branch
\[
  \lambda=\left(\frac{A_0(t)}w\right)^{1/D}
\]
near $\lambda=1$.  Since
\[
  \frac{P(t,r)}w=(\lambda\Phi_t(r))^D,
\]
\cref{lem:root-filter-general} with $m=2$ gives, on the tail,
\begin{align}\label{eq:root-tail-nonflat}
  \Bcal(w,t)
  &=\frac{\lambda^{D-2}}{DA_0(t)}
  \sum_{k=0}^{D-1}\omega^{-2k}
  \int_R^\infty
  e^{-r+\omega^k\lambda\Phi_t(r)}
  \frac{r}{\Phi_t(r)^2}\,dr
  +H_{\mathrm{corr}}(w,t),
\end{align}
where the finite correction $H_{\mathrm{corr}}$ is holomorphic near $(u_0,t_0)$.  Indeed, the root filter is used only on the radial tail, where $\lambda\Phi_t(r)$ stays away from zero; hence the negative powers occurring for $D=1,2$ have holomorphic, exponentially integrable coefficients.  For $k\neq0$, one has
\[
  \Re(1-\omega^k\lambda)>0
\]
uniformly after shrinking the neighborhood, so those tail integrals are holomorphic.  Only $k=0$ can be singular.

In the principal term put $Y=\Phi_t(r)$ and $r=\chi(t,Y)$.  Its amplitude becomes
\begin{equation}\label{eq:nonflat-amplitude}
  e^{Y-\chi(t,Y)}
  \frac{\chi(t,Y)}{Y^2}\,\partial_Y\chi(t,Y)
  =\frac{q_0(t)}Y+\frac{H(t,Y^{-1})}{Y^2},
\end{equation}
where
\[
  q_0(t)=\exp\!\left(\frac{A_1(t)}{DA_0(t)}\right)
\]
and $H$ is jointly holomorphic.  Apply Lemma~\ref{lem:nonflat-common-ray} to the
amplitude in \eqref{eq:nonflat-amplitude}.  Replacing the image curves
$\Phi_t([R,\infty))$ by one fixed ray changes the principal tail only by
a function jointly holomorphic in $(t,\eta)$ at $\eta=0$.

Put $\eta=1-\lambda$.  By \cref{lem:parameter-Laplace-log}, the principal tail is
\[
  \frac{\lambda^{D-2}}{DA_0(t)}
  \left[-q_0(t)\log\eta+\widetilde H_0(t,\eta)
  +\eta\log\eta\,\widetilde H_1(t,\eta)\right].
\]
Finally,
\[
  \eta=1-\left(\frac{A_0(t)}w\right)^{1/D}
  =\frac{w-A_0(t)}{DA_0(t)}
   +O\bigl((w-A_0(t))^2\bigr).
\]
Hence $\log\eta=\log(w-A_0(t))$ modulo a jointly holomorphic function, and $\eta\log\eta$ is $(w-A_0(t))\log(w-A_0(t))$ times a holomorphic unit, modulo a holomorphic function.  Evaluating the analytic prefactor at $\lambda=1$ gives \eqref{eq:boundary-symbol}.

  The lowest logarithmic coefficient is independent of the local
  $D$-th-root sheet.  Indeed, suppose a continued branch satisfies
  $\lambda\to\omega^q$ as $w\to A_0(t)$.  The unique potentially singular
  root-filter summand is then $k\equiv-q\pmod D$.  With
  $\mu=\omega^k\lambda\to1$, the same Laplace calculation applies with
  $\eta=1-\mu$, while its analytic prefactor has limiting root-of-unity
  factor
  \[
    \lambda^{D-2}\omega^{-2k}
    \longrightarrow
    \omega^{q(D-2)}\omega^{2q}=\omega^{qD}=1.
  \]
  Thus every continuation sheet has the coefficient displayed in
  \eqref{eq:boundary-symbol}.  This sheet-independence will be used in the
  constellation argument.
\end{proof}

\begin{lemma}[Continuation away from the leading divisor]
\label{lem:Bcal-continuation}
Let $U$ be a sufficiently small neighborhood on which $A_0$ is nowhere
zero and the uniform reversion of \cref{lem:uniform-nonflat-reversion} is
valid.  Put
\[
  \mathfrak X_U=
  \{(w,t)\in\C^\times\times U:w\ne A_0(t)\}.
\]
The germ $\Bcal(w,t)$ defined for large $w$ admits joint multivalued analytic
continuation along every path in $\mathfrak X_U$.  Every continued branch is
locally holomorphic on $\mathfrak X_U$.
\end{lemma}

\begin{proof}
Split the radial integral at a sufficiently large $R$.  On $0\le r\le R$,
\[
  \frac1w e^{-r}rE_{D,3}\!\left(\frac{P(t,r)}w\right)
\]
is jointly holomorphic for $w\ne0$, and compact integration preserves this
holomorphy.

For the tail, choose locally a branch
\[
  \lambda^D=\frac{A_0(t)}w.
\]
After the root-of-unity filter and the change of variable $Y=\Phi_t(r)$,
each exponential summand is, up to a locally holomorphic prefactor, of the
form
\begin{equation}\label{eq:continued-root-tail}
  L_k(t,\lambda)
  =\int_{\Gamma_t}e^{-\eta_kY}F(t,Y)\,dY,
  \qquad
  \eta_k=1-\omega^k\lambda,
\end{equation}
where
\[
  F(t,Y)=e^{Y-\chi(t,Y)}
  \frac{\chi(t,Y)}{Y^2}\partial_Y\chi(t,Y)=O(Y^{-1})
\]
is jointly holomorphic on a full common exterior neighborhood.

Fix a point $(t_*,\lambda_*)$ at which $\eta_k\ne0$.  Choose a ray
direction $\theta$ and a neighborhood $W$ of $(t_*,\lambda_*)$ such that
\[
  \Re(\eta_ke^{i\theta})\ge\delta>0
  \qquad\text{on }W.
\]
Inside the common exterior domain, join the initial point of $\Gamma_t$ to a
fixed sufficiently distant ray in direction $\theta$ by a finite connector,
chosen locally holomorphically in $t$.  The sum of the connector integral and
the integral along that ray converges normally and is jointly holomorphic on
$W$.

Whenever two such ray representations have a common decay region, Cauchy's
theorem on the truncated sector between the rays identifies them: the outer
circular arc tends to zero because of exponential decay and
$F(t,Y)=O(Y^{-1})$.  The admissible half-planes obtained by varying $\theta$
cover $\eta_k\ne0$, and along any path avoiding $0$ one may choose a finite
chain of overlapping half-planes.  Starting with the original decay region
of \eqref{eq:continued-root-tail}, the corresponding ray representations
therefore give its analytic continuation along the path.  Uniqueness on the
overlaps gives the asserted multivalued continuation.

A singularity of the $k$th summand can therefore occur only when
\[
  \eta_k=0
  \quad\Longleftrightarrow\quad
  \omega^k\lambda=1
  \quad\Longleftrightarrow\quad
  w=A_0(t).
\]
The algebraic parameter $\lambda$ can branch only around $w=0$, which is
already removed from $\mathfrak X_U$.

It remains to check the finite root-filter corrections.  For $D\ge3$ there
are none.  For $D=2$ the correction reduces on the tail to an exponentially
weighted multiple of $1/P(t,r)$, and for $D=1$ to a combination of
$1/P(t,r)$ and $w/P(t,r)^2$.  After increasing $R$, the polynomial $P(t,r)$
is uniformly nonzero there, so these correction integrals are jointly
holomorphic.  This proves the asserted continuation and local holomorphy.
\end{proof}

\begin{remark}[Why the lower layers do not enter the symbol]\label{rem:exact-symbol}
The degree-$(D-1)$ layer changes the constant term in the reverted coordinate and exponentiates to the factor in \eqref{eq:boundary-symbol}.  The layers $A_2,\ldots,A_D$ first occur in the $Y^{-2}$ and lower terms of \eqref{eq:nonflat-amplitude}; their Laplace transforms therefore begin at order $(w-A_0)\log(w-A_0)$.
\end{remark}

\section{Positive-interval constellations}\label{sec:constellation}

Let
\[
  A=A_0\in\C[t],\qquad d=\deg A\ge1.
\]
Choose a sufficiently large regular value $c$ of $A$.  Join $c$ by pairwise disjoint arcs to every finite critical value of $A$, and also to $A(0)$ and $A(1)$ if either is not already a critical value.  Let $S$ be the resulting star.  The preimage $A^{-1}(S)$ is a planar tree, and there is a unique oriented path $\mu_{0,1}$ from $0$ to $1$.

Let
\[
  t_1(u),\ldots,t_d(u)
\]
be the inverse branches of $A$ near $c$.  For each arm $\gamma_s$ of the star, the path $\mu_{0,1}$ determines an incidence vector
\[
  f_s=(f_{s,1},\ldots,f_{s,d})\in\Z^d.
\]
Let $G_A$ be the monodromy group and define
\begin{equation}\label{eq:incidence-module}
  M_{A,0,1}=\Span_\Q\{\sigma f_s:\sigma\in G_A,\ s\}.
\end{equation}
This is the module $M_{P,a,b}$ of Pakovich and Muzychuk \cite[\S2.1]{PakovichMuzychuk}.

Number the branches at infinity so that monodromy around infinity is $(1\ 2\ \cdots\ d)$, and put
\[
  \omega=e^{2\pi i/d},
  \qquad
  \mathbf w_j=(1,\omega^j,\omega^{2j},\ldots,\omega^{(d-1)j}).
\]
For a rational subspace $M\subseteq\Q^d$, write
\[
  M_\C^\perp
  =\left\{v\in\C^d:\sum_{i=1}^d m_iv_i=0
    \text{ for every }m=(m_i)\in M\right\}.
\]
Thus $M_\C^\perp$ is taken with respect to the complex-bilinear extension
of the rational dot product, not a Hermitian form.

\begin{proposition}[Primitive Fourier inclusion; Pakovich--Muzychuk]\label{prop:PM-Fourier}
If $d\ge2$, then $(M_{A,0,1})_\C$ contains $\mathbf w_1$ and $\mathbf w_{d-1}$.
\end{proposition}

\begin{proof}
Pakovich and Muzychuk prove that $M_{A,0,1}$ contains a distinguished irreducible rational monodromy subspace $W$ \cite[Proposition~4.1 and its proof]{PakovichMuzychuk}.  Their proof explicitly places $\mathbf w_1$ in $W_\C$.  Since $W$ is defined over $\Q$, it is stable under complex conjugation, so $\mathbf w_{d-1}=\overline{\mathbf w_1}$ also belongs to $W_\C$.
\end{proof}

For $|w|$ large, $\Bcal(w,t)\,dt$ is holomorphic in $t$.  Deforming $[0,1]$ to $\mu_{0,1}$ and changing variables $u=A(t)$ on each lifted edge gives
\begin{equation}\label{eq:constellation-decomposition}
  \Kcal(w)=\sum_s\int_{\gamma_s}
  \sum_{i=1}^d f_{s,i}
  \frac{\Bcal(w,t_i(u))}{A'(t_i(u))}\,du.
\end{equation}
This is the same geometric decomposition as
\cite[(16)--(17)]{PakovichMuzychuk}; only the branch density differs.  The
edge integrals are understood as improper integrals at critical-value
vertices.  If $u_*$ is such a vertex, a lifted branch tends to a point
$\tau$ of ramification index $e$, and locally
\[
  t_i(u)-\tau=O\bigl((u-u_*)^{1/e}\bigr),
  \qquad
  \frac1{A'(t_i(u))}=\frac{dt_i}{du}
  =O\bigl(|u-u_*|^{1/e-1}\bigr).
\]
Since $1/e-1>-1$, these endpoint singularities are integrable.  In
particular, the deformation and every later edgewise integration are
well-defined.

Define
\begin{equation}\label{eq:E-and-v}
  E(t)=\exp\!\left(\frac{A_1(t)}{DA(t)}\right),
  \qquad
  v_i(u)=\frac{E(t_i(u))}{A'(t_i(u))}.
\end{equation}
For an arm $\gamma_s$, put
\begin{equation}\label{eq:arm-density}
  F_s(w,u)=\sum_{i=1}^d f_{s,i}
  \frac{\Bcal(w,t_i(u))}{A'(t_i(u))}.
\end{equation}

\begin{lemma}[Terminal continuation and arm isolation]
\label{lem:arm-continuation}
Choose the terminal subarcs of the star pairwise disjoint except at their
common regular endpoint $c$, and choose their interiors to avoid $0$ and a
prescribed finite set of points.  After shortening a terminal subarc
$\tau_s$ if necessary, the following hold on every continuation sheet
reached from large $w$.

First, in a slit neighborhood of the diagonal $w=u$ over $u\in\tau_s$, the
density \eqref{eq:arm-density} has the joint expansion
\begin{equation}\label{eq:tubular-arm-expansion}
  F_s(w,u)
  =g_{s,0}(u)\log(w-u)
   +(w-u)\log(w-u)G_s(w,u)+H_s(w,u),
\end{equation}
where $G_s$ and $H_s$ are jointly holomorphic and
\begin{equation}\label{eq:g-s0}
  g_{s,0}(u)=-\frac1{Du}\sum_{i=1}^d f_{s,i}v_i(u).
\end{equation}
Second, for every interior point $w_0\in\tau_s\setminus\{c\}$ sufficiently
close to $c$, there is a disk $D_{w_0}$ whose closure meets no other arm and
no prescribed point.  On $D_{w_0}$, the sum of all other arm integrals,
together with the part of the $s$-th arm outside $\tau_s$, is single-valued
and holomorphic.  Thus only the integral over $\tau_s$ can have a lateral
jump across the selected arm at $w_0$.
\end{lemma}

\begin{proof}
The terminal subarc contains no critical value, so the branches $t_i(u)$ and
the factors $A'(t_i(u))^{-1}$ are holomorphic on a neighborhood of
$\tau_s$.  The parameterized boundary symbol, including the
sheet-independence proved at the end of \cref{thm:boundary-symbol}, gives
\eqref{eq:tubular-arm-expansion} jointly in $(w,u)$; summing its lowest
coefficients gives \eqref{eq:g-s0}.  Uniformity on the compact shortened
subarc permits termwise integration.

Fix $w_0\in\tau_s\setminus\{c\}$ sufficiently close to $c$.  Since the
terminal arms are disjoint away from $c$, a small closed disk around $w_0$
meets no other arm, no prescribed point, and no part of the selected arm
outside $\tau_s$.  For every integration point on those omitted pieces one
has $w\ne u$ throughout the disk.  By
\cref{lem:Bcal-continuation}, every already continued branch of the
integrand is locally holomorphic there.  The endpoint estimates preceding
\eqref{eq:E-and-v} give a locally uniform integrable majorant at
critical-value vertices, so integration preserves holomorphy.  The disk is
simply connected, hence these local branches are single-valued on it.  This
proves the asserted isolation of the selected terminal integral.
\end{proof}

\begin{lemma}[Terminal lateral jump]\label{lem:Volterra}
Let $\tau=\gamma([0,1])$ be an oriented analytic arc, where $\gamma$ extends
holomorphically to a neighborhood of $[0,1]$, $\gamma'(1)\neq0$, and
$\gamma(1)=c$.  Suppose that on a slit tube around the terminal part of
$\tau$ one has
\[
  F(w,u)=g_0(u)\log(w-u)
  +(w-u)\log(w-u)G(w,u)+H(w,u),
\]
with $g_0$, $G$, and $H$ holomorphic and with $H$ single-valued.  Let
$I(w)=\int_\tau F(w,u)\,du$, and let $I_+$ and $I_-$ denote its two lateral
continuations across the terminal subarc.  For $w=\gamma(s)$ with $s<1$
sufficiently close to $1$,
\begin{equation}\label{eq:terminal-jump-formula}
\begin{aligned}
 I_+(\gamma(s))-I_-(\gamma(s))
  ={}&2\pi i\,\varepsilon_\tau
  \int_s^1\Bigl[g_0(\gamma(v))\\
 &\quad +(\gamma(s)-\gamma(v))
  G(\gamma(s),\gamma(v))\Bigr]\gamma'(v)\,dv.
\end{aligned}
\end{equation}
where $\varepsilon_\tau\in\{\pm1\}$ depends only on the orientation and the
choice of lateral sides.  In particular, if the lateral jump vanishes,
then $g_0(c)=0$.
\end{lemma}

\begin{proof}
Across the chosen slit, $\log(w-u)$ changes by the constant
$2\pi i\varepsilon_\tau$ precisely for the portion of the oriented arc
between $w$ and the terminal endpoint $c$; it has no jump on the remaining
portion.  This proves \eqref{eq:terminal-jump-formula}.

Let the integral on the right, without its nonzero constant factor, be
$J(s)$.  Differentiating with respect to $s$ gives
\[
  J'(s)=-g_0(\gamma(s))\gamma'(s)+R(s),
\]
where $R(s)\to0$ as $s\to1$: every term in $R(s)$ is an integral over
$[s,1]$ of a bounded holomorphic function, and the lower-end contribution
from the factor $\gamma(s)-\gamma(v)$ is zero.  If the jump vanishes, then
$J\equiv0$, so
\[
  0=\lim_{s\to1}J'(s)=-g_0(c)\gamma'(1).
\]
Since $\gamma'(1)\neq0$, one obtains $g_0(c)=0$.
\end{proof}

\begin{proposition}[Rationality forces incidence orthogonality]\label{prop:rational-orthogonality}
Assume that $\Kcal$ is rational.  For a sufficiently large regular center
$c$, chosen away from $0$ and the poles of $\Kcal$,
\begin{equation}\label{eq:incidence-relations}
  \sum_{i=1}^d f_{s,i}v_i(c)=0
  \qquad\text{for every arm }s.
\end{equation}
After analytic continuation, the same identity holds for every monodromy
conjugate of every $f_s$.  Equivalently,
\begin{equation}\label{eq:v-orthogonal}
  v(c)\in(M_{A,0,1})_\C^\perp.
\end{equation}
\end{proposition}

\begin{proof}
Choose $c$ outside the finite critical-value set, outside $0$, and outside
the finite pole set of the assumed rational function $\Kcal$.  Choose the
interiors of the terminal subarcs to avoid the latter two finite sets.  On
the overlap with the original large-$w$ domain, the continued constellation
sum equals the rational function $\Kcal$.

Fix an arm $s$ and a point $w_0$ in its terminal subarc, sufficiently close
to $c$ but distinct from $c$.  Choose the disk $D_{w_0}$ supplied by
\cref{lem:arm-continuation}, small enough to contain no pole of $\Kcal$.
The two lateral continuations of the total constellation sum across the
selected arm agree on this disk because $\Kcal$ is single-valued.  Every
other contribution has zero jump there, so the lateral jump of the isolated
terminal integral vanishes.  This holds for every such $w_0$.  Applying
\cref{lem:Volterra} to \eqref{eq:tubular-arm-expansion} gives
$g_{s,0}(c)=0$, and \eqref{eq:g-s0} yields
\eqref{eq:incidence-relations}.

It remains to justify the monodromy closure.  Choose a small simply
connected disk $\Omega$ of regular values around $c$, disjoint from $0$ and
from the poles of $\Kcal$.  The inverse branches $t_i(u)$ and the vectors
$v(u)$ are holomorphic on $\Omega$.  Move the common endpoint of the star
inside $\Omega$, keeping the remote parts of the arms fixed and identifying
the terminal lifts by the local covering trivialization.  The incidence
vectors are then locally constant.  Repeating the preceding lateral-jump
argument for every $c'\in\Omega$ proves
\[
  \sum_i f_{s,i}v_i(c')=0
  \qquad(c'\in\Omega),
\]
so this is an identity of analytic branches, not merely an equality at one
point.

Continue that identity around loops in the regular-value locus while
avoiding $0$ and the finite pole set of $\Kcal$.  These additional punctures
are regular values: every ordinary monodromy loop can be perturbed to avoid
them, and a small loop around one of them has trivial covering monodromy.
Thus all elements of $G_A$ are still realized.  Upon return, the inverse
branches are permuted by the corresponding monodromy element $\sigma$.
With the chosen left-action convention the continued relation is
$(\sigma f_s)\cdot v=0$; with the opposite convention it is
$(\sigma^{-1}f_s)\cdot v=0$, which generates the same module.  Thus every
monodromy conjugate is orthogonal to $v$, which is exactly
\eqref{eq:v-orthogonal} with the complex-bilinear convention fixed above.
\end{proof}

If $d\ge2$, combining
\cref{prop:rational-orthogonality,prop:PM-Fourier} yields, on a sufficiently
small regular-value disk $\Omega$ around $c$, the identity of analytic
branches
\begin{equation}\label{eq:primitive-relation}
  \sum_{i=1}^d\omega^{-(i-1)}v_i(u)=0
  \qquad(u\in\Omega).
\end{equation}

\section{The infinity obstruction and the diagonal theorem}\label{sec:infinity}

\begin{lemma}[Rotated exponential averages do not vanish]\label{lem:rotated-exp}
Let
\[
  T(U)=\alpha U+\beta+O(U^{-1}),
  \qquad \alpha\neq0,
\]
be holomorphic for large $U$, let $R\in\C(t)$ be rational, and let $\omega=e^{2\pi i/d}$.  Then
\[
  \sum_{j=0}^{d-1}e^{R(T(\omega^jU))}T'(\omega^jU)
\]
is not identically zero near infinity.
\end{lemma}

\begin{proof}
Write $R(T(U))=\Phi(U)+O(U^{-1})$, where $\Phi$ is the
polynomial part of the Laurent expansion.  Partition the indices into
classes for which the polynomials $\Phi(\omega^jU)$ are identical.
Every rotated phase has the same constant term $\Phi(0)$; consequently,
two distinct rotated phase polynomials cannot differ by a nonzero
constant.  Their difference, when nonzero, has positive degree.  Choose
a ray avoiding the finitely many anti-Stokes directions determined by
these nonzero differences.  Along that ray one class $C$ has strictly
maximal real part for large $U$.  Its contribution is
\[
  |C|\alpha e^{\Phi_C(U)}(1+o(1)),
\]
while every other class is exponentially smaller.  This is nonzero.  If all phase polynomials agree, the sum is $d\alpha e^{\Phi(U)}(1+o(1))$, again nonzero.
\end{proof}

\begin{theorem}[Diagonal positive-interval constellation theorem]\label{thm:diagonal}
Assume $A_0$ is nonconstant.  Then
\[
  \Jcal(s)=\sum_{n\ge0}\frac{I_n}{Dn+2}s^n
\]
is not rational.  Consequently, $I_n\neq0$ for infinitely many $n$.
\end{theorem}

\begin{proof}
Suppose $\Jcal$ were rational.  Then $\Kcal(w)=w^{-1}\Jcal(w^{-1})$ would be rational.  Put $A=A_0$ and $d=\deg A$.

If $d=1$, the lifted path from $0$ to $1$ is nontrivial, so some arm has $f_{s,1}\neq0$.  Proposition \ref{prop:rational-orthogonality} would then force the nonzero density $E(t_1(c))/A'(t_1(c))$ to vanish, a contradiction.

Assume $d\ge2$.  Choose a Laurent inverse at infinity
\begin{equation}\label{eq:T-inverse}
  A(T(U))=U^d,
  \qquad
  T(U)=\alpha U+\beta+O(U^{-1}),
  \quad\alpha\neq0.
\end{equation}
Choose $U_0$ large with $c=U_0^d$, and number the branches compatibly so
that, on a connected exterior sector containing $U_0$,
\[
  t_i(U^d)=T(\omega^{i-1}U).
\]
Put
\[
  R(t)=\frac{A_1(t)}{DA(t)}.
\]
Differentiating \eqref{eq:T-inverse} gives
\[
  A'(T(V))T'(V)=dV^{d-1}.
\]
Hence, after pulling the analytic identity
\eqref{eq:primitive-relation} back under $u=U^d$,
\[
  0=\sum_{i=1}^d\omega^{-(i-1)}v_i(U^d)
   =\frac1{dU^{d-1}}
    \sum_{i=1}^d
    e^{R(T(\omega^{i-1}U))}T'(\omega^{i-1}U)
\]
for $U$ in a neighborhood of $U_0$.  The last sum is holomorphic on the
connected exterior sector.  By the identity theorem it vanishes identically
there, contradicting \cref{lem:rotated-exp}.  Hence $\Jcal$ is not rational.
If only finitely many $I_n$ were nonzero, then $\Jcal$ would be a polynomial,
which is impossible.
\end{proof}


\begin{theorem}[Nonflat-top exclusion]\label{thm:nonflat-exclusion}
Let $0\neq f\in\C[x,y]$ have exact degree $D$.  If
$A_0(t)=f_D(t,1-t)$ is nonconstant, then $\Lfac(f^n)\neq0$ for infinitely
many $n\ge1$.  In particular, $f$ is not a counterexample to
\cref{thm:FC2}.
\end{theorem}

\begin{proof}
Apply \cref{thm:diagonal} to the projective layers of $f$.  By
\cref{prop:radial-projective},
\[
  I_n=\frac{\Lfac(f^n)}{\Gamma(Dn+2)}.
\]
Thus infinitely many factorial moments are nonzero.
\end{proof}

\begin{corollary}[Flat-top reduction]\label{cor:flat-top-reduction}
Every hypothetical nonzero counterexample to \cref{thm:FC2} has
\[
  f_D(x,y)=a(x+y)^D
\]
for some $a\in\C^\times$.
\end{corollary}

\begin{proof}
By \cref{thm:nonflat-exclusion}, $f_D(t,1-t)$ is constant.  The final
assertion of \cref{lem:projective-degree} gives the stated form.
\end{proof}


\section{The first flat logarithmic symbol}\label{sec:flat-symbol}

Assume that $f$ has exact degree $D\ge1$ and, after multiplication by a
nonzero scalar,
\begin{equation}\label{eq:flat-normalized}
  f_D=(x+y)^D.
\end{equation}
Then
\begin{equation}\label{eq:flat-P}
  P(t,r)=f(rt,r(1-t))
  =r^D+A_1(t)r^{D-1}+\cdots+A_D(t).
\end{equation}
Define
\[
  F(t,u)=1+A_1(t)u+\cdots+A_D(t)u^D.
\]
Choose a relatively compact complex neighborhood $U$ of $[0,1]$ and
$\rho>0$ so small that $F(t,u)$ is nonzero for $t\in U$ and $|u|<\rho$.
On this product take the unique holomorphic branch
\[
  \Psi(t,u)=F(t,u)^{1/D},
  \qquad \Psi(t,0)=1,
\]
and put
\begin{equation}\label{eq:flat-phi}
  \phi_t(r)=r\Psi(t,r^{-1}).
\end{equation}
Thus $P(t,r)=\phi_t(r)^D$.

\begin{lemma}[Uniform flat radial inverse]\label{lem:flat-inverse}
After shrinking $U$ if necessary, there is a common exterior neighborhood
on which $\phi_t$ has a holomorphic inverse $r=\chi(t,Y)$, jointly
holomorphic in $(t,Y)$, satisfying
\begin{equation}\label{eq:flat-chi}
  \chi(t,Y)=Y-\frac{A_1(t)}D+O(Y^{-1}).
\end{equation}
Consequently,
\begin{equation}\label{eq:flat-amplitude-leading}
  e^{Y-\chi(t,Y)}\frac{\chi(t,Y)}Y\partial_Y\chi(t,Y)
  =\exp\!\left(\frac{A_1(t)}D\right)+O(Y^{-1}),
\end{equation}
locally uniformly in $t\in U$.
\end{lemma}

\begin{proof}
Locally uniformly for $t\in U$,
\[
  \phi_t(r)=r+\frac{A_1(t)}D+O(r^{-1}).
\]
In the coordinate $z=r^{-1}$, the map
$z\mapsto1/\phi_t(1/z)$ has derivative $1$ at $z=0$.  The
parameter-dependent inverse-function theorem therefore gives a common
inverse at infinity and \eqref{eq:flat-chi}.  Substitution into the left
side of \eqref{eq:flat-amplitude-leading} gives the displayed constant
term.
\end{proof}

\begin{lemma}[Common-ray comparison]\label{lem:flat-common-ray}
After shrinking $U$ and increasing $R$, there are a common exterior sector
$\mathfrak S$ and a ray $\ell=Y_0+[0,\infty)\subset\mathfrak S$ such that
all curves $\Gamma_t=\phi_t([R,\infty))$, $t\in U$, lie in $\mathfrak S$
and have bounded displacement from $\ell$.  Let $F(t,Y)$ be jointly
holomorphic on $U\times\mathfrak S$, analytic in $Y^{-1}$ at infinity, and
satisfy
\[
  F(t,Y)=q_0(t)+O(Y^{-1})
\]
locally uniformly in $t$.  For $\eta$ in a sufficiently small decay sector,
the difference
\[
  \int_{\Gamma_t}e^{-\eta Y}F(t,Y)\,dY
  -\int_\ell e^{-\eta Y}F(t,Y)\,dY
\]
extends jointly holomorphically across $\eta=0$.
\end{lemma}

\begin{proof}
The expansion of $\phi_t$ shows that the curves, the common ray, and the
initial and outer connectors used below all lie in one common exterior
sector.  Write $F=q_0+\widetilde F$ with
$\widetilde F=O(Y^{-1})$.  For the constant term, path independence of a
primitive of $e^{-\eta Y}$ gives
\[
  q_0(t)\frac{e^{-\eta a_t}-e^{-\eta Y_0}}{\eta},
  \qquad a_t=\phi_t(R),
\]
which is holomorphic at $\eta=0$.  For $\widetilde F$, truncate the two
curves at a common large ray parameter.  Their outer endpoints remain at
bounded distance, so the outer connector has bounded length and contributes
$O(|Y|^{-1})$, uniformly on closed decay subsectors; this tends to zero.
Cauchy's theorem, now applied inside the domain of holomorphy
$U\times\mathfrak S$, reduces the difference to a finite initial-connector
integral, which is jointly holomorphic.
\end{proof}

Define
\begin{equation}\label{eq:flat-generating}
  \Gcal_f(s)=\sum_{n\ge0}
  \frac{\Lfac(f^n)}{\Gamma(Dn+2)}s^n.
\end{equation}

\begin{proposition}[First flat singular coefficient]\label{prop:first-flat-symbol}
Let $\lambda^D=s$ near $\lambda=1$ and put $\eta=1-\lambda$.  On a fixed
admissible sector at $\eta=0$,
\begin{equation}\label{eq:first-flat-symbol}
  \Gcal_f(\lambda^D)
  =\frac{c_0(f)}{D\lambda\eta}+O(\log\eta),
\end{equation}
where the error is the sum of a holomorphic germ and a function bounded by
$C(1+|\log\eta|)$ on every smaller closed sector, and
\begin{equation}\label{eq:flat-c0}
  c_0(f)=\int_0^1
  \exp\!\left(\frac{A_1(t)}D\right)\,dt.
\end{equation}
In particular,
\begin{equation}\label{eq:moments-imply-c0}
  \Lfac(f^n)=0\quad(n\ge1)
  \qquad\Longrightarrow\qquad c_0(f)=0.
\end{equation}
\end{proposition}

\begin{proof}
By \cref{prop:Mittag-representations},
\[
  \Gcal_f(s)=\int_0^1\int_0^\infty
  e^{-r}rE_{D,2}(sP(t,r))\,dr\,dt.
\]
Split the radial integral at a sufficiently large $R$ and apply
\cref{lem:root-filter-general} with $m=1$ on the tail, where $\phi_t(r)$ is
uniformly nonzero.  The compact radial part is holomorphic near
$\lambda=1$.  The low-degree correction vanishes for $D\ge2$; when $D=1$,
$C_{1,1}(X)=X^{-1}$ and its tail contribution is
\[
  -\frac1\lambda\int_0^1\int_R^\infty
  \frac{e^{-r}}{\Psi(t,r^{-1})}\,dr\,dt,
\]
which is also holomorphic near $\lambda=1$.  Every root-filter term with
$k\neq0$ is holomorphic there because
$\Re(1-\omega^k\lambda)>0$.  The only singular term is
\begin{equation}\label{eq:flat-principal-tail}
  \frac1{D\lambda}\int_0^1\int_R^\infty
  \frac{e^{-r+\lambda\phi_t(r)}}{\Psi(t,r^{-1})}\,dr\,dt.
\end{equation}
Set $Y=\phi_t(r)$ and $r=\chi(t,Y)$.  Since
$\Psi(t,r^{-1})=Y/r$, \eqref{eq:flat-principal-tail} becomes
\[
  \frac1{D\lambda}\int_0^1\int_{\Gamma_t}
  e^{-\eta Y}
  e^{Y-\chi(t,Y)}\frac{\chi(t,Y)}Y\partial_Y\chi(t,Y)
  \,dY\,dt.
\]
By \cref{lem:flat-inverse,lem:flat-common-ray}, modulo a function
holomorphic at $\eta=0$ this equals
\[
  \frac1{D\lambda}\int_{Y_0}^\infty e^{-\eta Y}
  \left[c_0(f)+O(Y^{-1})\right]dY.
\]
The constant term contributes
$c_0(f)e^{-\eta Y_0}/(D\lambda\eta)$, whose difference from the displayed
pole in \eqref{eq:first-flat-symbol} is holomorphic.  The Laplace transform
of the $O(Y^{-1})$ remainder is $O(1+|\log\eta|)$ on smaller closed sectors.
This proves \eqref{eq:first-flat-symbol} and \eqref{eq:flat-c0}.

For real $0<\lambda<1$, the root-filter integrals above converge absolutely
and agree algebraically with the original Mittag--Leffler integral.  Hence
the sectorial germ just computed is the analytic continuation, along the
positive interval, of the original germ of $\Gcal_f(\lambda^D)$ at
$\lambda=0$.

If all positive factorial moments vanish, then that original germ is
$\Gcal_f(s)\equiv1$, and uniqueness of analytic continuation makes the
continued germ constant as well.  A constant function has no pole at
$\lambda=1$, so \eqref{eq:first-flat-symbol} forces $c_0(f)=0$.
\end{proof}


\section{Algebraic exponential integrals}\label{sec:exp-integrals}

\begin{theorem}[Algebraic exponential-integral theorem]\label{thm:exp-integral}
Let $q\in\Qbar[x]$ be nonconstant, and let $A,B\in\Qbar$ with $A\ne B$.
Then
\[
  \int_A^B e^{q(x)}\,dx
\]
is transcendental.  In particular, it is nonzero.
\end{theorem}

The integral is independent of the path because the integrand is entire.
The rest of the paper proves this theorem before returning to
\cref{thm:FC2}.

\section{\texorpdfstring{$E$}{E}-functions and algebraic normalization}\label{sec:E-functions}

We use the following definition, in the form adopted by Beukers \cite[\S1]{Beukers}.

\begin{definition}[$E$-function]\label{def:E}
An entire function
\[
  F(z)=\sum_{m\ge0} a_m\frac{z^m}{m!}
\]
is an $E$-function if:
\begin{enumerate}[label=\textup{(\roman*)},leftmargin=2.2em]
\item $a_m\in\Qbar$ for every $m$;
\item the logarithmic height $h(a_0,\ldots,a_m)$ is $O(m)$;
\item $F$ satisfies a nonzero linear differential equation with coefficients in $\Qbar[z]$.
\end{enumerate}
\end{definition}

The key arithmetic input is Beukers' lifting theorem.

\begin{theorem}[Beukers]\label{thm:Beukers}
Let $F_1,\ldots,F_N$ be $E$-functions forming a vector solution of
\[
  \mathbf F'=A(z)\mathbf F,
  \qquad A(z)\in M_N(\Qbar(z)),
\]
and let $T(z)\in\Qbar[z]$ be a common denominator of the entries of $A(z)$.
Let $\xi\in\Qbar$ satisfy $\xi T(\xi)\ne0$.  If
$P\in\Qbar[X_1,\ldots,X_N]$ is homogeneous and
\[
  P\bigl(F_1(\xi),\ldots,F_N(\xi)\bigr)=0,
\]
then there is a polynomial
$Q\in\Qbar[z,X_1,\ldots,X_N]$, homogeneous in the variables
$X_1,\ldots,X_N$ and of the same degree as $P$, such that
\[
  Q\bigl(z,F_1(z),\ldots,F_N(z)\bigr)\equiv0,
  \qquad
  Q(\xi,X_1,\ldots,X_N)=P(X_1,\ldots,X_N).
\]
In particular, a linear relation at $z=\xi$ lifts to a linear functional
relation whose coefficient polynomials specialize coefficientwise to the
prescribed relation.  No hypothesis of $\C(z)$-linear independence of the
$F_i$ is required.
\end{theorem}

This is \cite[Theorem~1.3]{Beukers}; its linear-relation form is
\cite[Theorem~3.2]{Beukers}.

We shall also use the elementary independence of distinct exponentials.

\begin{lemma}[Distinct exponentials]\label{lem:exp-independent}
If $\lambda_1,\ldots,\lambda_s\in\C$ are distinct, then
\[
  e^{\lambda_1 z},\ldots,e^{\lambda_s z}
\]
are linearly independent over $\C(z)$.
\end{lemma}

\begin{proof}
Suppose a nontrivial relation with the fewest nonzero terms exists.  Differentiate it and eliminate one term.  Minimality then forces the logarithmic derivative of a nonzero rational function to be a nonzero constant.  This is impossible, since a logarithmic derivative has no nonzero constant term at infinity.
\end{proof}

\subsection{Algebraic normalization}

Let $q_0$ have degree $n\ge2$.  Translation by an algebraic number removes the coefficient of $x^{n-1}$, and an algebraic nonzero scaling of the variable normalizes the leading coefficient to $1/n$.  Thus there are $\alpha\in\Qbar^\times$ and $\beta\in\Qbar$ such that
\begin{equation}\label{eq:normalized-q}
  q(y):=q_0(\alpha y+\beta)
  =\frac{y^n}{n}+a_{n-2}y^{n-2}+\cdots+a_1y+a_0,
  \qquad a_j\in\Qbar.
\end{equation}
The change of variable $x=\alpha y+\beta$ gives
\[
  \int_A^B e^{q_0(x)}\,dx
  =\alpha\int_{(A-\beta)/\alpha}^{(B-\beta)/\alpha}
  e^{q(y)}\,dy.
\]
Hence the transformed endpoints remain distinct and algebraic, and the original integral differs from the normalized one by a nonzero algebraic factor.  Since such multiplication preserves transcendence, it suffices to prove \cref{thm:exp-integral} under \eqref{eq:normalized-q}.  We impose this normalization through \cref{sec:exp-main-proof}.

\section{The intrinsic moment connection}\label{sec:exp-moment}

Assume \eqref{eq:normalized-q}, put $n=\deg q$, and set
\[
  K=\Qbar(z).
\]
Define the twisted de Rham quotient
\begin{equation}\label{eq:twisted-module}
  \Mcal_q
  =K[x]\big/(\partial_x+zq'(x))K[x].
\end{equation}
We write $[p]$ for the class of $p\in K[x]$.

\begin{lemma}[Brieskorn basis]\label{lem:Brieskorn-basis}
The classes
\[
  e_j=[x^j],
  \qquad 0\le j\le n-2,
\]
form a $K$-basis of $\Mcal_q$.  More precisely,
\[
  \mathscr G_q
  =\C[z,z^{-1},x]\Big/
  (\partial_x+zq'(x))\C[z,z^{-1},x]
\]
is a free $\C[z,z^{-1}]$-module with the same basis.  Differentiation in $z$ defines the connection
\begin{equation}\label{eq:brieskorn-connection}
  \nabla_{\partial_z}[p]=[\partial_zp+q p].
\end{equation}
\end{lemma}

\begin{proof}
The relation
\[
  [zq'(x)p(x)]=-[p'(x)]
\]
reduces every monomial of degree at least $n-1$ to lower degree, both over $K$ and over $\C[z,z^{-1}]$.  Thus the displayed classes span.

If $0\ne p\in K[x]$, then $(\partial_x+zq')p$ has $x$-degree $\deg_x p+n-1$, and its leading term cannot cancel.  The same argument applies over $\C[z,z^{-1}]$.  Hence the image contains no nonzero polynomial of degree at most $n-2$, proving independence.

The connection is well defined because
\[
  [\partial_z+q,\,\partial_x+zq']=0:
\]
indeed, the commutators $[\partial_z,zq']=q'$ and $[q,\partial_x]=-q'$ cancel.  Formula \eqref{eq:brieskorn-connection} is the algebraic form of differentiating the kernel $e^{zq}$.
\end{proof}

For $0\le j\le n-2$, perform Euclidean division
\begin{equation}\label{eq:division}
  q(x)x^j=h_j(x)q'(x)+r_j(x),
  \qquad \deg r_j\le n-2.
\end{equation}
Let $\mathsf C$ and $\mathsf D$ be the matrices whose $j$th columns are the coefficient vectors of $r_j$ and $-h_j'$ in the basis $1,x,\ldots,x^{n-2}$.

\begin{lemma}[The matrices $\mathsf C$ and $\mathsf D$]\label{lem:CD}
The following assertions hold.
\begin{enumerate}[label=\textup{(\roman*)},leftmargin=2.2em]
\item $\mathsf C$ is multiplication by $q$ on the Jacobian algebra $\Qbar[x]/(q')$;
\item $\mathsf D$ is triangular with diagonal
\[
  -\frac1n,-\frac2n,\ldots,-\frac{n-1}{n};
\]
\item the connection matrix on $\Mcal_q$ is $\mathsf C+\mathsf D/z$.
\end{enumerate}
\end{lemma}

\begin{proof}
In $\Mcal_q$,
\[
  [h_jq']=-\frac1z[h_j'],
\]
so \eqref{eq:division} gives
\[
  [qx^j]=[r_j]-\frac1z[h_j'].
\]
This proves (i) and (iii).

Because $q$ has no $x^{n-1}$ term and $q'$ has no $x^{n-2}$ term, division gives
\[
  h_j(x)=\frac{x^{j+1}}n+
  \text{a polynomial of degree at most }j-1.
\]
Therefore
\[
  -h_j'(x)=-\frac{j+1}{n}x^j+
  \text{terms of degree at most }j-2,
\]
which proves (ii).
\end{proof}

The degree of $q$ is denoted by $n$ throughout the exponential-integral
argument.  In \cref{sec:exp-moment,sec:exp-semisimple,sec:exp-rigidity,sec:exp-relations},
the capital letter $D$ introduced below denotes a residue matrix and is
unrelated to the bivariate degree used in \cref{sec:radial,sec:flat-symbol}.

Let $A\ne B$ be algebraic, and define
\begin{equation}\label{eq:Uj}
  U_j(z)=\int_A^B x^j e^{zq(x)}\,dx,
  \qquad
  \mathbf U=(U_0,\ldots,U_{n-2})^{\mathsf T}.
\end{equation}
Define the endpoint vector
\begin{equation}\label{eq:b-x}
  \mathbf b(x)=n\bigl(h_0(x),\ldots,h_{n-2}(x)\bigr)^{\mathsf T}.
\end{equation}
Since $h_0(x)=x/n$, the first coordinate of $\mathbf b(x)$ is $x$.

Let
\[
  \Lambda_{\partial}=\{q(A),q(B)\}
\]
with repetitions removed, and put
\begin{equation}\label{eq:b-lambda}
  \mathbf b_\lambda
  =\sum_{\substack{x\in\{A,B\}\\q(x)=\lambda}}
  \varepsilon_x\mathbf b(x),
  \qquad
  \varepsilon_B=1,
  \quad \varepsilon_A=-1.
\end{equation}

\begin{proposition}[Moment system]\label{prop:moment-system}
Put $C=\mathsf C^{\mathsf T}$ and $D=\mathsf D^{\mathsf T}$.  Then
\begin{equation}\label{eq:moment-system}
  \mathbf U'
  =\left(C+\frac Dz\right)\mathbf U
  +\frac1{nz}\sum_{\lambda\in\Lambda_{\partial}}
  \mathbf b_\lambda e^{\lambda z}.
\end{equation}
Moreover,
\begin{equation}\label{eq:D-spectrum}
  \Spec(D)=
  \left\{-\frac1n,-\frac2n,\ldots,-\frac{n-1}{n}\right\},
\end{equation}
and
\begin{equation}\label{eq:block-sum}
  \sum_{\lambda\in\Lambda_{\partial}}
  (\mathbf b_\lambda)_0=B-A\ne0.
\end{equation}
\end{proposition}

\begin{proof}
From \eqref{eq:division},
\begin{align*}
  U_j'
  &=\int_A^B q(x)x^j e^{zq(x)}\,dx\\
  &=\int_A^B\left(r_j(x)-\frac1z h_j'(x)\right)e^{zq(x)}\,dx
    +\frac1z\left[h_j(x)e^{zq(x)}\right]_{A}^{B}.
\end{align*}
Writing these identities simultaneously gives \eqref{eq:moment-system}.  The spectral assertion follows from \cref{lem:CD}.  The first coordinate of $\mathbf b(x)$ is $x$, which gives \eqref{eq:block-sum}.
\end{proof}


\begin{lemma}[Joint height from one common denominator]\label{lem:joint-height}
Let $K$ be a number field, let $\alpha_0,\ldots,\alpha_m\in K$, and let
$\delta\in\Z_{\ge1}$ satisfy $\delta\alpha_k\in\mathcal O_K$ for every
$k$.  Then
\[
  h(\alpha_0,\ldots,\alpha_m)
  \le \log\delta
  +\log\max\!\left(
    1,\max_{0\le k\le m}\max_{\sigma:K\hookrightarrow\C}
    |\sigma(\alpha_k)|
  \right).
\]
Here the left side is the absolute logarithmic Weil height of
$(1:\alpha_0:\cdots:\alpha_m)$.
\end{lemma}

\begin{proof}
With the standard absolute values and local degrees $n_v=[K_v:\Q_v]$,
\[
  h(\alpha_0,\ldots,\alpha_m)
  =\frac1{[K:\Q]}\sum_{v\in M_K}n_v
   \log\max\bigl(1,|\alpha_0|_v,\ldots,|\alpha_m|_v\bigr).
\]
Put $\beta_k=\delta\alpha_k\in\mathcal O_K$.  Projective invariance gives
\[
  (1:\alpha_0:\cdots:\alpha_m)
  =(\delta:\beta_0:\cdots:\beta_m).
\]
At every finite place, all displayed coordinates have absolute value at most
$1$, so the corresponding logarithm is at most zero.  At an infinite place,
\[
  \max\bigl(|\delta|_v,|\beta_0|_v,\ldots,|\beta_m|_v\bigr)
  \le
  \delta\max\!\left(
    1,\max_{k,\sigma}|\sigma(\alpha_k)|
  \right).
\]
Since $\sum_{v\mid\infty}n_v=[K:\Q]$, summing these upper bounds proves the
claim.
\end{proof}

\begin{lemma}[Moment functions are $E$-functions]\label{lem:E-functions}
For $0\le j\le n-2$, the entire function $U_j$ in \eqref{eq:Uj} is an
$E$-function.  More precisely, if
\[
  U_j(z)=\sum_{m\ge0}u_{j,m}\frac{z^m}{m!},
  \qquad
  u_{j,m}=\int_A^B x^j q(x)^m\,dx,
\]
then there is a constant $\kappa_j$ such that
\[
  h(u_{j,0},\ldots,u_{j,m})\le \kappa_j(m+1)
  \qquad(m\ge0),
\]
and $U_j$ satisfies a nonzero homogeneous linear differential equation over
$\Qbar[z]$ of order at most $n+1$.
\end{lemma}

\begin{proof}
We integrate over the straight segment from $A$ to $B$; the value is path
independent because the integrand is entire in $x$.  Uniform convergence on
compact subsets gives
\[
  U_j(z)=\sum_{m\ge0}u_{j,m}\frac{z^m}{m!},
  \qquad
  u_{j,m}=\int_A^B x^j q(x)^m\,dx,
\]
and also proves that $U_j$ is entire.

Let
\[
  K_0=\Q\bigl(A,B,\text{the coefficients of }q\bigr).
\]
Write
\[
  q(x)=\sum_{i=0}^n c_i x^i,
  \qquad
  q(x)^m=\sum_{k=0}^{nm}\gamma_{m,k}x^k.
\]
For $m\ge1$,
\[
  \gamma_{m,k}
  =\sum_{i_1+\cdots+i_m=k}c_{i_1}\cdots c_{i_m},
\]
while $\gamma_{0,0}=1$.  Termwise integration gives the exact formula
\begin{equation}\label{eq:u-exact}
  u_{j,m}
  =\sum_{k=0}^{nm}\gamma_{m,k}
   \frac{B^{j+k+1}-A^{j+k+1}}{j+k+1}.
\end{equation}
It follows immediately that $u_{j,m}\in K_0$.

We next bound all archimedean conjugates simultaneously.  Put
\[
  M=\max_{\sigma:K_0\hookrightarrow\C}
     \max\bigl(1,|\sigma(A)|,|\sigma(B)|\bigr),
\]
\[
  \Theta=\max_{\sigma:K_0\hookrightarrow\C}
          \max_{|x|\le M}|q^\sigma(x)|,
  \qquad
  \Delta=\max_{\sigma:K_0\hookrightarrow\C}
          |\sigma(B)-\sigma(A)|.
\]
Every conjugate segment lies in $|x|\le M$, and applying $\sigma$ to
\eqref{eq:u-exact} gives the corresponding conjugate integral.  Hence
\[
  |\sigma(u_{j,m})|
  \le \Delta M^j\Theta^m.
\]
With $C_1=\max(1,\Delta,M,\Theta)$, this is at most
$C_1^{m+j+1}$ for every embedding $\sigma$.

Choose $d\in\Z_{\ge1}$ so that
\[
  dc_0,\ldots,dc_n,dA,dB\in\mathcal O_{K_0}.
\]
For $m\ge0$, set
\[
  \delta_m
  =d^{(n+1)m+j+1}
   \operatorname{lcm}(1,\ldots,nm+j+1).
\]
Then $d^m\gamma_{m,k}$ and
$d^{j+k+1}(B^{j+k+1}-A^{j+k+1})$ are algebraic integers.  Since
$0\le k\le nm$, the $k$th summand of \eqref{eq:u-exact}, after multiplication
by $\delta_m$, is
\[
  \bigl(d^m\gamma_{m,k}\bigr)d^{nm-k}
  \Bigl(d^{j+k+1}(B^{j+k+1}-A^{j+k+1})\Bigr)
  \frac{\operatorname{lcm}(1,\ldots,nm+j+1)}{j+k+1},
\]
a product of algebraic integers.  Thus $\delta_m u_{j,m}$ is integral.
Moreover $\delta_\ell\mid\delta_m$ whenever $0\le\ell\le m$, so $\delta_m$
is a common denominator for $u_{j,0},\ldots,u_{j,m}$.

By \cref{lem:joint-height},
\[
  h(u_{j,0},\ldots,u_{j,m})
  \le \log\delta_m+(m+j+1)\log C_1.
\]
The elementary Chebyshev bound
$\log\operatorname{lcm}(1,\ldots,N)\le2N$ gives
\[
  \log\delta_m
  \le \bigl((n+1)m+j+1\bigr)\log d+2(nm+j+1).
\]
Since $j$ is fixed, these inequalities give
$h(u_{j,0},\ldots,u_{j,m})=O(m)$.

It remains to establish a scalar differential equation.  Order the distinct
endpoint values as $\lambda_1,\ldots,\lambda_s$, where $s\le2$, and let
$\mathcal W$ be the $\Qbar(z)$-span of
\[
  U_0,\ldots,U_{n-2},
  e^{\lambda_1z},\ldots,e^{\lambda_sz}
\]
inside the field of meromorphic functions on $\C\setminus\{0\}$.  By
\cref{prop:moment-system}, this space is stable under $\partial_z$, and
\[
  \dim_{\Qbar(z)}\mathcal W\le n-1+s\le n+1.
\]
Therefore the $n+2$ functions
$U_j,\partial_zU_j,\ldots,\partial_z^{n+1}U_j$ are linearly dependent over
$\Qbar(z)$.  Clearing denominators gives a nonzero operator in
$\Qbar[z]\langle\partial_z\rangle$ of order at most $n+1$ annihilating
$U_j$.  The three defining conditions of an $E$-function now follow.
\end{proof}

\section{Semisimplicity of the moment connection}\label{sec:exp-semisimple}

Let $\Sigma$ be the finite set of branch values of $q$, and set
\[
  V=\Aone_t\setminus\Sigma.
\]
Over $V$, the map $q:q^{-1}(V)\to V$ is a degree-$n$ unramified covering.  Let $\Lcal$ be its permutation local system and
\[
  \Lcal^0=\ker(\Tr:\Lcal\to\C_V)
\]
its reduced, or augmentation, local system of rank $n-1$.

\begin{proposition}[Semisimple Fourier--Laplace connection]\label{prop:semisimple}
After extension of constants to $\C$, the differential module over $\C(z)$
defined by the homogeneous system
\begin{equation}\label{eq:homogeneous-system}
  \mathbf Y'=\left(C+\frac Dz\right)\mathbf Y
\end{equation}
is semisimple.
\end{proposition}

\begin{proof}
The cover $q^{-1}(V)\to V$ is connected, because $q^{-1}(V)$ is obtained
from the affine line by deleting finitely many points.  Its monodromy action
on a fiber is therefore transitive.  The permutation local system has the
trace splitting
\[
  \Lcal\simeq\C_V\oplus\Lcal^0:
\]
indeed, the diagonal map followed by the trace is multiplication by
$n$, which is invertible over $\C$.  The monodromy image is finite, so
Maschke's theorem makes $\Lcal^0$ semisimple.  Write
\[
  \Lcal^0=\bigoplus_\alpha\Lcal_\alpha
\]
as a direct sum of irreducible local systems.

A finite morphism has no positive-dimensional fibers, so the defining
condition for smallness is vacuous; hence $q$ is small.  The proper small-map
theorem gives
\[
  Rq_*\C_{\Aone_x}[1]\simeq\operatorname{IC}(\Lcal);
\]
see \cite[Proposition~4.2.1 and Remark~4.2.4]{deCataldoMigliorini}.  The
intermediate extensions $\operatorname{IC}(\Lcal_\alpha)$ are simple
perverse sheaves, while
$\operatorname{IC}(\C_V)=\C_{\Aone_t}[1]$.  Under the Riemann--Hilbert
correspondence they correspond to simple regular holonomic
$\Dmod_{\Aone_t}$-modules; see
\cite[\S\S3.4, 7.1--7.2, and 8.2.2]{HTT}.  Since the direct image above is
perverse, the derived $D$-module direct image is concentrated in degree zero,
and compatibility with proper direct image yields
\begin{equation}\label{eq:direct-image-splitting}
  q_+\Ocal_{\Aone_x}
  \simeq
  \Ocal_{\Aone_t}\oplus\Ncal^0,
\end{equation}
where $\Ncal^0$ is a direct sum of simple regular holonomic modules.

Use a Fourier variable $\tau$ and the Weyl-algebra convention
\begin{equation}\label{eq:Fourier-convention}
  t\longmapsto-\partial_\tau,
  \qquad
  \partial_t\longmapsto\tau.
\end{equation}
It corresponds to the kernel $e^{-\tau t}$.  This is an isomorphism of Weyl
algebras, so Fourier transform is an exact equivalence and preserves simple
objects and direct sums.  The transform of the constant connection is
\[
  \widehat{\Ocal_{\Aone_t}}
  \simeq
  \Dmod_{\Aone_\tau}/\Dmod_{\Aone_\tau}\tau,
\]
which is supported at $\tau=0$ and disappears after localization to
$\C(\tau)$.

We record why generic localization preserves each surviving simple summand.
Let $M$ be a simple Weyl module and let $S=\C[\tau]\setminus\{0\}$.  Its
$S$-torsion is a Weyl submodule: if $p(\tau)m=0$, then
$p(\tau)^2\partial_\tau m=0$.  Thus either $S^{-1}M=0$, or
$M\hookrightarrow S^{-1}M$.  In the latter case, a nonzero differential
submodule $N\subseteq S^{-1}M$ meets $M$ nontrivially after clearing a
denominator.  The intersection $N\cap M$ is a nonzero Weyl submodule, hence
is $M$, and therefore $N=S^{-1}M$.  Thus every nonzero localized simple
summand is irreducible.

It remains to identify the localized Fourier transform and its signs.  Factor
$q$ as its graph embedding followed by the projection to the $t$-line.  The
graph module is generated by a symbol $\delta_q$ satisfying
\[
  (t-q(x))\delta_q=0,
  \qquad
  (\partial_x+q'(x)\partial_t)\delta_q=0.
\]
After \eqref{eq:Fourier-convention}, the transformed generator satisfies
\[
  (\partial_\tau+q(x))\widehat\delta_q=0,
  \qquad
  (\partial_x+\tau q'(x))\widehat\delta_q=0;
\]
it is the exponential generator $e^{-\tau q(x)}$.  On its coefficient
polynomials, the relative $x$-differential and the $\tau$-connection are
therefore
\[
  \partial_x-\tau q'(x),
  \qquad
  \partial_\tau-q(x),
\]
respectively.  The projection along the $x$-line is computed by the relative
de Rham complex
\[
  \left[
    \C(\tau)[x]
    \xrightarrow{\ \partial_x-\tau q'\ }
    \C(\tau)[x]
  \right],
\]
with the left term in degree $-1$ and the right term in degree $0$.  Its
negative-degree cohomology is the kernel, which is zero: for nonzero
$p\in\C(\tau)[x]$, the term $-\tau q'p$ has strictly larger $x$-degree than
$p'$.  Consequently the degree-zero generic Fourier transform is the
cokernel
\begin{equation}\label{eq:negative-kernel-module}
  \C(\tau)[x]\Big/(\partial_x-\tau q'(x))\C(\tau)[x],
  \qquad
  \nabla_{\partial_\tau}[p]=[\partial_\tau p-q p].
\end{equation}

Pulling back by $\tau=-z$ changes the relative operator to
$\partial_x+zq'$ and, because $\partial_z=-\partial_\tau$, changes the
connection to $\partial_z+q$.  Thus \eqref{eq:negative-kernel-module} becomes
$\Mcal_q$ with the connection \eqref{eq:brieskorn-connection}.  By
\cref{lem:CD}, its matrix in the monomial basis is
$\mathsf C+\mathsf D/z$.  It follows from the preceding decomposition and
localization argument that
$\Mcal_q\otimes_{\Qbar(z)}\C(z)$ is semisimple.

For clarity, if a connection has matrix $M$ in a basis $(e_j)$, so that
$\nabla e_j=\sum_iM_{ij}e_i$, a horizontal functional has value column
$\mathbf y=(\phi(e_j))_j$ satisfying $\mathbf y'=M^{\mathsf T}\mathbf y$.
Thus the homogeneous period system has coefficient matrix
\[
  \mathsf C^{\mathsf T}+\frac{\mathsf D^{\mathsf T}}z
  =C+\frac Dz.
\]
The row module with connection
$\boldsymbol\phi\mapsto\boldsymbol\phi'+\boldsymbol\phi(C+D/z)$ used below
is its dual.  Duality preserves semisimplicity, proving the proposition.
\end{proof}

\section{Horizontal-endomorphism rigidity}\label{sec:exp-rigidity}

We now prove the main structural input.  The proof is entirely algebraic and treats degenerate critical points and coincident critical values uniformly.

Let
\[
  \Vcal=\C[x]_{<n-1}.
\]
For $g\in\Vcal$, divide
\begin{equation}\label{eq:g-division}
  gq=h_gq'+r_g,
  \qquad \deg r_g<n-1,
\end{equation}
and define linear operators
\begin{equation}\label{eq:CstarDstar}
  \mathsf C_*g=r_g,
  \qquad
  \mathsf D_*g=-h_g'.
\end{equation}
The quotient and remainder depend linearly on $g$ by uniqueness of Euclidean
division.  Moreover,
\[
  \deg h_g\le\deg g+1\le n-1,
  \qquad
  \deg(-h_g')\le n-2,
\]
so both operators map $\Vcal$ to itself.  They are represented by
$\mathsf C$ and $\mathsf D$ in the monomial basis.

For a critical point $\tau$, put
\[
  m_\tau=\ord_\tau q',
  \qquad
  \lambda_\tau=q(\tau).
\]

\begin{lemma}[Critical-point basis and local spectrum]\label{lem:critical-basis}
For every critical point $\tau$ and $1\le k\le m_\tau$, define
\begin{equation}\label{eq:g-tau-k}
  g_{\tau,k}(x)=\frac{q'(x)}{(x-\tau)^k}.
\end{equation}
Then:
\begin{enumerate}[label=\textup{(\roman*)},leftmargin=2.2em]
\item the polynomials $g_{\tau,k}$, over all $\tau$ and $k$, form a basis of
$\Vcal$;
\item $\mathsf C_*$ is semisimple and
\[
  \mathsf C_*g_{\tau,k}=\lambda_\tau g_{\tau,k};
\]
\item put
\[
  E_\tau=\Span_\C\{g_{\tau,1},\ldots,g_{\tau,m_\tau}\}.
\]
If $E_\lambda=\ker(\mathsf C_*-\lambda I)$ and $\pi_\lambda$ is the spectral
projection, then
\[
  E_\lambda=\bigoplus_{\tau:\lambda_\tau=\lambda}E_\tau,
  \qquad
  \overline D_\lambda=\pi_\lambda\mathsf D_*|_{E_\lambda}
\]
preserves every summand $E_\tau$.  In the ordered basis
$g_{\tau,1},\ldots,g_{\tau,m_\tau}$, its restriction to $E_\tau$ is
triangular with diagonal
\begin{equation}\label{eq:local-D-spectrum}
  -\frac{m_\tau}{m_\tau+1},
  -\frac{m_\tau-1}{m_\tau+1},
  \ldots,
  -\frac1{m_\tau+1}.
\end{equation}
\end{enumerate}
\end{lemma}

\begin{proof}
Factor
\[
  q'(x)=c\prod_\tau(x-\tau)^{m_\tau}.
\]
Reduction modulo $q'$ identifies $\Vcal$ with the Jacobian algebra, and the
Chinese remainder theorem gives
\[
  \C[x]/(q')
  \simeq
  \bigoplus_\tau\C[x]/(x-\tau)^{m_\tau}.
\]
The $E_\sigma$-component of an element of $\Vcal$ is precisely its residue
class modulo $(x-\sigma)^{m_\sigma}$.  For fixed $\tau$, the images of
$g_{\tau,k}$ vanish in every other local factor, while in the $\tau$-factor,
with $y=x-\tau$, they are
\[
  g_{\tau,k}=u_\tau(y)y^{m_\tau-k},
  \qquad u_\tau(0)\ne0.
\]
Hence $g_{\tau,1},\ldots,g_{\tau,m_\tau}$ form a basis of the $\tau$-factor.
Since $\sum_\tau m_\tau=n-1=\dim\Vcal$, all the displayed vectors form a
basis of $\Vcal$, proving (i).

Since $q-\lambda_\tau$ vanishes at $\tau$ to order $m_\tau+1$, the polynomial
\[
  h_{\tau,k}(x)=\frac{q(x)-\lambda_\tau}{(x-\tau)^k}
\]
is well defined and satisfies
\[
  qg_{\tau,k}=h_{\tau,k}q'+\lambda_\tau g_{\tau,k}.
\]
Thus $\mathsf C_*g_{\tau,k}=\lambda_\tau g_{\tau,k}$.  The basis just
constructed is therefore an eigenbasis for $\mathsf C_*$, proving (ii), and
\[
  E_\lambda=\bigoplus_{\tau:\lambda_\tau=\lambda}E_\tau.
\]

Write $m=m_\tau$, $y=x-\tau$, and
\[
  q(x)-\lambda_\tau=ay^{m+1}+O(y^{m+2}),
  \qquad a\ne0.
\]
Then
\[
  g_{\tau,k}=a(m+1)y^{m-k}+O(y^{m-k+1})
\]
and
\[
  -h_{\tau,k}'
  =-a(m+1-k)y^{m-k}+O(y^{m-k+1}).
\]
In the ordered local basis, the error term can involve only
$g_{\tau,j}$ with $j<k$: it has no local term of order below $y^{m-k}$.
Consequently the compressed local block is triangular and its $k$th diagonal
entry is
\[
  -\frac{m+1-k}{m+1}.
\]

It remains to exclude cross terms between distinct critical points over the
same critical value.  If $\sigma\ne\tau$ and
$\lambda_\sigma=\lambda_\tau$, then
\[
  h_{\tau,k}=\frac{q-\lambda_\tau}{(x-\tau)^k}
\]
vanishes at $\sigma$ to order $m_\sigma+1$, because the denominator is a unit
there.  Hence $h_{\tau,k}'$ vanishes modulo
$(x-\sigma)^{m_\sigma}$.  Thus the compression
$\pi_\lambda\mathsf D_*|_{E_\lambda}$ preserves each $E_\tau$ and has the
claimed diagonal spectrum.
\end{proof}

\begin{remark}\label{rem:full-vs-local-D}
There is no conflict between \eqref{eq:D-spectrum} and
\eqref{eq:local-D-spectrum}.  The first is the spectrum of the full operator
$\mathsf D_*$, whereas the second describes its block-diagonal compression
relative to the spectral decomposition of $\mathsf C_*$.  Their traces agree:
\[
  -\sum_\tau\sum_{r=1}^{m_\tau}\frac r{m_\tau+1}
  =-\frac12\sum_\tau m_\tau
  =-\frac{n-1}{2}
  =-\sum_{j=1}^{n-1}\frac jn.
\]
\end{remark}

\begin{theorem}[Horizontal-endomorphism rigidity]\label{thm:end-rigidity}
Let $P(z)\in M_{n-1}(\C(z))$ satisfy
\begin{equation}\label{eq:endomorphism-equation}
  P'=\left[C+\frac Dz,P\right].
\end{equation}
Then $P$ is constant.  Consequently, the horizontal endomorphism matrices of
the homogeneous moment system \eqref{eq:homogeneous-system} are exactly
\[
  Z(C)\cap Z(D),
\]
and all are constant.  Equivalently, in the monomial basis the horizontal
endomorphism matrices of the algebraic connection $\Mcal_q$ are exactly
\[
  Z(\mathsf C)\cap Z(\mathsf D).
\]
The same constancy conclusion holds for the corresponding dual modules.
\end{theorem}

\begin{proof}
First exclude finite poles.  If $P$ had a pole of order $m\ge1$ at
$z_0\ne0$, then $P'$ would have pole order $m+1$, whereas the right side of
\eqref{eq:endomorphism-equation} would have order at most $m$.  Thus no such
pole occurs.

At $z=0$, suppose
\[
  P=z^{-m}P_{-m}+O(z^{-m+1}),
  \qquad m\ge1,
  \quad P_{-m}\ne0.
\]
The coefficient of $z^{-m-1}$ is
\[
  -mP_{-m}=[D,P_{-m}].
\]
The eigenvalues of $D$ are $-j/n$, $1\le j\le n-1$, so the spectrum of
$\operatorname{ad}D:X\mapsto[D,X]$ consists of their pairwise differences
and is contained strictly in $(-1,1)$.  It does not contain the negative
integer $-m$, a contradiction.  Hence $P$ has no finite pole.  Since $P$ is
rational, it is a polynomial matrix.

We now exclude positive degree.  Transpose
\eqref{eq:endomorphism-equation} and put $Q=P^{\mathsf T}$.  Then
\begin{equation}\label{eq:row-endomorphism}
  Q'=[Q,\mathsf C_*+\mathsf D_*/z].
\end{equation}
Suppose
\[
  Q(z)=Q_dz^d+Q_{d-1}z^{d-1}+\cdots,
  \qquad d\ge1,
  \quad Q_d\ne0.
\]
Multiplying \eqref{eq:row-endomorphism} by $z$ and comparing the coefficients
of $z^{d+1}$ and $z^d$ gives
\begin{align}
  [Q_d,\mathsf C_*]&=0,\label{eq:leading-commute}\\
  dQ_d&=[Q_{d-1},\mathsf C_*]+[Q_d,\mathsf D_*].
  \label{eq:next-leading}
\end{align}
This part of the argument uses only the matrix semisimplicity of
$\mathsf C_*$ proved in \cref{lem:critical-basis}, not the differential-module
semisimplicity of \cref{prop:semisimple}.

For each critical value $\lambda$, its spectral projection is a polynomial in
$\mathsf C_*$:
\[
  \pi_\lambda
  =\prod_{\mu\ne\lambda}
    \frac{\mathsf C_*-\mu I}{\lambda-\mu}.
\]
Equation \eqref{eq:leading-commute} therefore implies that $Q_d$ commutes with
every $\pi_\lambda$ and is block diagonal.  Put
\[
  Q_{d,\lambda}=\pi_\lambda Q_d\pi_\lambda.
\]
For at least one $\lambda$, this block is nonzero.  Applying $\pi_\lambda$ on
both sides of \eqref{eq:next-leading} kills the
$[Q_{d-1},\mathsf C_*]$ term and gives
\[
  \pi_\lambda[Q_d,\mathsf D_*]\pi_\lambda
  =[Q_{d,\lambda},\pi_\lambda\mathsf D_*\pi_\lambda]
  =[Q_{d,\lambda},\overline D_\lambda].
\]
Thus
\begin{equation}\label{eq:ad-local-D}
  dQ_{d,\lambda}
  =[Q_{d,\lambda},\overline D_\lambda].
\end{equation}

After triangularizing $\overline D_\lambda$, the commutator operator
$X\mapsto[X,\overline D_\lambda]$ is triangular on matrix units, with diagonal
entries equal to pairwise differences of eigenvalues of
$\overline D_\lambda$.  By \eqref{eq:local-D-spectrum}, all those eigenvalues
lie strictly in $(-1,0)$, so every difference lies strictly in $(-1,1)$.
Equation \eqref{eq:ad-local-D} would make the positive integer $d$ an
eigenvalue of this commutator operator, which is impossible.  Hence $d=0$.

Therefore $P$ is constant.  Substitution into
\eqref{eq:endomorphism-equation} gives $[C,P]=[D,P]=0$, and every constant
matrix commuting with $C$ and $D$ is plainly horizontal.  Transposition gives
the statement for $\Mcal_q$.  If a dual endomorphism acts on row vectors by
right multiplication, the calculation
\[
  \nabla(\boldsymbol\phi P)=(\nabla\boldsymbol\phi)P
  \quad\Longleftrightarrow\quad
  P'=[C+D/z,P]
\]
shows that the same constancy result applies to the dual modules.
\end{proof}

\begin{corollary}[Multiplicity-free semisimplicity]\label{cor:multiplicity-free}
The horizontal endomorphism algebra of the moment connection is a
commutative finite-dimensional $\C$-algebra.  In particular, the semisimple
moment connection is multiplicity-free: its simple constituents are
pairwise nonisomorphic.
\end{corollary}

\begin{proof}
The matrix $D$ has distinct eigenvalues, so it is diagonalizable and its
centralizer is commutative.  Apply \cref{thm:end-rigidity}.  For a semisimple
differential module, a repeated simple constituent contributes a full
matrix algebra to the endomorphism algebra; commutativity therefore excludes
repetitions.
\end{proof}

\begin{remark}
The multiplicity-free conclusion is also visible geometrically.  Monodromy
at infinity is an $n$-cycle, so its cyclic subgroup acts simply transitively
on a generic fiber; the restriction of the permutation representation to
that subgroup is its regular representation and is multiplicity-free.  If an
irreducible constituent of the full monodromy representation occurred twice,
then every cyclic character appearing in its restriction would occur at least
twice, a contradiction.
\end{remark}

\section{Relation modules and constant projectors}\label{sec:exp-relations}

Let
\[
  \Lambda=\{0,q(A),q(B)\}
\]
with repetitions removed, and put
\begin{equation}\label{eq:E-space}
  \Ecal=
  \Span_{\C(z)}\{e^{\lambda z}:\lambda\in\Lambda\}.
\end{equation}
If $0$ is not an endpoint value, its forcing block is understood to be zero.

The functions
\[
  U_0,\ldots,U_{n-2},
  \{e^{\lambda z}:\lambda\in\Lambda\}
\]
form a vector of $E$-functions: the endpoint exponents $\lambda$ are
algebraic, and $e^{\lambda z}$ is an $E$-function.  Adjoining these exponential
coordinates to \eqref{eq:moment-system} gives a homogeneous block system over
$\Qbar(z)$ with common denominator $z$.  When $0$ is not an endpoint value,
the coordinate $e^{0z}=1$ is adjoined with a zero forcing column.  This
coordinate will homogenize the value relation used below.

Define the relation module over $\C(z)$ by
\begin{equation}\label{eq:relation-module}
  \Rcal=
  \{\boldsymbol\phi\in\C(z)^{1\times(n-1)}:
  \boldsymbol\phi\mathbf U\in\Ecal\}.
\end{equation}

\begin{lemma}[Relation subconnection]\label{lem:relation-subconnection}
The space $\Rcal$ is a differential submodule of the dual homogeneous connection:
\[
  \boldsymbol\phi\in\Rcal
  \quad\Longrightarrow\quad
  \boldsymbol\phi'
  +\boldsymbol\phi\left(C+\frac Dz\right)
  \in\Rcal.
\]
\end{lemma}

\begin{proof}
Let
\[
  A(z)=C+\frac Dz,
  \qquad
  \mathbf f(z)=\frac1{nz}\sum_{\lambda\in\Lambda}
  \mathbf b_\lambda e^{\lambda z},
\]
where $\mathbf b_0=0$ when necessary.  From \eqref{eq:moment-system},
\[
  \left(\boldsymbol\phi\mathbf U\right)'
  =\left(\boldsymbol\phi'+\boldsymbol\phi A\right)\mathbf U
   +\boldsymbol\phi\mathbf f.
\]
The left side belongs to $\Ecal$ because $\Ecal$ is stable under differentiation, and $\boldsymbol\phi\mathbf f\in\Ecal$.  Hence
\[
  \left(\boldsymbol\phi'+\boldsymbol\phi A\right)\mathbf U\in\Ecal,
\]
which is the required assertion.
\end{proof}

\begin{proposition}[A value relation forces an exponential polynomial]\label{prop:value-to-exp-poly}
If $U_0(1)$ is algebraic, then
\begin{equation}\label{eq:exp-polynomial-U0}
  U_0(z)=\sum_{\lambda\in\Lambda}p_\lambda(z)e^{\lambda z}
\end{equation}
for polynomials $p_\lambda\in\C[z]$.
\end{proposition}

\begin{proof}
Let $a=U_0(1)\in\Qbar$.  Order
\[
  \Lambda=\{0,\lambda_2,\ldots,\lambda_s\}
\]
and form the full column
\[
  \mathbf F(z)
  =\bigl(
     U_0(z),\ldots,U_{n-2}(z),
     1,e^{\lambda_2z},\ldots,e^{\lambda_sz}
   \bigr)^{\mathsf T}.
\]
Put
\[
  B_\Lambda=\frac1n
  \begin{pmatrix}
    \mathbf b_0&\mathbf b_{\lambda_2}&\cdots&\mathbf b_{\lambda_s}
  \end{pmatrix},
  \qquad
  \Delta_\Lambda=\diag(0,\lambda_2,\ldots,\lambda_s),
\]
where $\mathbf b_0=0$ if $0\notin\Lambda_\partial$.  Then
\[
  \mathbf F'
  =\begin{pmatrix}
      C+D/z&B_\Lambda/z\\
      0&\Delta_\Lambda
    \end{pmatrix}\mathbf F.
\]
The constant coordinate $1=e^{0z}$ makes
$U_0(1)-a\cdot1=0$ a homogeneous linear relation.  More explicitly, let
\[
  \boldsymbol\alpha
  =\bigl(1,0,\ldots,0;\,-a,0,\ldots,0\bigr),
\]
where the semicolon separates the moment and exponential coordinates.  Then
$\boldsymbol\alpha\mathbf F(1)=0$.  The block-system matrix has common
denominator $T(z)=z$, and the point $\xi=1$ satisfies
$\xi T(\xi)=1\ne0$.  By \cref{thm:Beukers}, this prescribed linear form lifts
to a polynomial coefficient row
\[
  R(z)\in\Qbar[z]^{1\times((n-1)+s)}
\]
such that
\[
  R(z)\mathbf F(z)=0,
  \qquad
  R(1)=\boldsymbol\alpha.
\]
Writing
\[
  R(z)=\bigl(
    \boldsymbol\phi(z);
    c_0(z),c_{\lambda_2}(z),\ldots,c_{\lambda_s}(z)
  \bigr)
\]
gives
\[
  \boldsymbol\phi(z)\mathbf U(z)
  +c_0(z)+\sum_{j=2}^s c_{\lambda_j}(z)e^{\lambda_jz}=0,
  \qquad
  \boldsymbol\phi(1)=\mathbf e_0:=(1,0,\ldots,0).
\]
Because $0\in\Lambda$, the entire right-hand exponential combination belongs
to $\Ecal$; hence $\boldsymbol\phi\in\Rcal$.  We now regard this algebraic
relation inside the complex differential module $\Rcal\subseteq\C(z)^{n-1}$.

By \cref{prop:semisimple}, the homogeneous moment module and its dual are
semisimple over $\C(z)$.  The differential submodule $\Rcal$ therefore has a
horizontal complement and a horizontal idempotent projector.  In row
coordinates the projector acts by right multiplication by a rational matrix
$P(z)$ with image $\Rcal$.  Horizontality is
\[
  P'=\left[C+\frac Dz,P\right].
\]
By \cref{thm:end-rigidity}, $P$ is constant; substituting into the equation
also gives $[C,P]=[D,P]=0$.  Let
\[
  W=\operatorname{im}_{\C}(\,\cdot P)
  \subseteq\C^{1\times(n-1)}.
\]
Then $\Rcal=W\otimes_\C\C(z)$.  Since
$\boldsymbol\phi=\boldsymbol\phi P$ and $P$ is constant, evaluation at
$z=1$ gives
\[
  \mathbf e_0=\boldsymbol\phi(1)=\boldsymbol\phi(1)P\in W.
\]
Thus the constant row $\mathbf e_0$ lies in $\Rcal$, and
$U_0=\mathbf e_0\mathbf U\in\Ecal$.

It remains to upgrade rational exponential coefficients to polynomial ones.
Choose a constant full-row-rank matrix $S$ whose rows form a basis of $W$, and
put $\mathbf V=S\mathbf U$.  The space $W$ is invariant under right
multiplication by $C$ and $D$: if $w=uP$, then
\[
  wC=uPC=uCP=(uC)P\in W,
  \qquad
  wD=uPD=uDP=(uD)P\in W.
\]
Hence there are unique constant matrices $C_W,D_W$ such that
\[
  SC=C_WS,
  \qquad
  SD=D_WS.
\]
With $\mathbf b_{W,\lambda}=S\mathbf b_\lambda$, where
$\mathbf b_0=0$ when $0\notin\Lambda_\partial$, the moment system restricts
to
\begin{equation}\label{eq:restricted-system}
  \mathbf V'
  =\left(C_W+\frac{D_W}{z}\right)\mathbf V
  +\frac1{nz}\sum_{\lambda\in\Lambda}
   \mathbf b_{W,\lambda}e^{\lambda z}.
\end{equation}
Every row of $S$ belongs to $\Rcal$, so every component of $\mathbf V$ lies
in $\Ecal$.  By \cref{lem:exp-independent}, there are unique rational vectors
$\mathbf r_\lambda(z)\in\C(z)^{\dim W}$ such that
\[
  \mathbf V(z)=\sum_{\lambda\in\Lambda}
  \mathbf r_\lambda(z)e^{\lambda z}.
\]
Equating the independent exponential coefficients in
\eqref{eq:restricted-system} gives
\begin{equation}\label{eq:restricted-block}
  \mathbf r_\lambda'
  =(C_W-\lambda I)\mathbf r_\lambda
  +\frac{D_W}{z}\mathbf r_\lambda
  +\frac1{nz}\mathbf b_{W,\lambda}.
\end{equation}

A pole of $\mathbf r_\lambda$ at $z_0\ne0$ is impossible: differentiation
raises its order by one, whereas all terms on the right have at most the
original order.  At $z=0$, suppose
\[
  \mathbf r_\lambda=z^{-m}v+O(z^{-m+1}),
  \qquad m\ge1,
  \quad v\ne0.
\]
The forcing term has pole order at most $1$, strictly below $m+1$ even when
$m=1$.  Comparing the coefficient of $z^{-m-1}$ therefore gives
\[
  -mv=D_Wv.
\]
Every eigenvalue of $D_W$ is an eigenvalue of $D$.  Indeed, if
$u^{\mathsf T}D_W=\mu u^{\mathsf T}$, then full row rank of $S$ gives
$u^{\mathsf T}S\ne0$, and
\[
  (u^{\mathsf T}S)D
  =u^{\mathsf T}D_WS
  =\mu(u^{\mathsf T}S).
\]
By \eqref{eq:D-spectrum}, $\Spec(D_W)\subset(-1,0)$, so it contains no
negative integer $-m$.  This contradiction excludes a pole at zero.  Thus
every rational vector $\mathbf r_\lambda$ has no finite pole and belongs to
$\C[z]^{\dim W}$.

Finally choose a constant row $a_0$ with $\mathbf e_0=a_0S$.  Then
\[
  U_0=\mathbf e_0\mathbf U=a_0\mathbf V
  =\sum_{\lambda\in\Lambda}
   \bigl(a_0\mathbf r_\lambda(z)\bigr)e^{\lambda z},
\]
and $p_\lambda=a_0\mathbf r_\lambda\in\C[z]$.  In fact, the argument proves
that every component of $\mathbf V$ is an exponential polynomial.
\end{proof}

\begin{remark}\label{rem:entire-rational-coefficients}
Entirety alone does not imply polynomial exponential coefficients:
\[
  \frac{e^z-1}{z}
  =\frac1z e^z-\frac1z
\]
is entire although both displayed coefficients have a pole.  This is
$U_0(z)$ for $q(x)=x$, $A=0$, $B=1$; that linear case has rank-zero moment
connection and is handled separately by Lindemann--Weierstrass.  For
$n\ge2$, the residue nonresonance in \eqref{eq:D-spectrum} is therefore
load-bearing.
\end{remark}

\section{Constant-coefficient annihilation}\label{sec:exp-annihilation}

Assume \eqref{eq:exp-polynomial-U0}.  Choose an integer $N$ larger than the degrees of all nonzero polynomials $p_\lambda$; equivalently, with the convention $\deg0=-\infty$, choose
\[
  N>\max_{\lambda\in\Lambda}\deg p_\lambda.
\]
Define the constant-coefficient differential operator
\begin{equation}\label{eq:L-operator}
  L=\prod_{\lambda\in\Lambda}(\partial_z-\lambda)^N.
\end{equation}
Then $LU_0=0$.  Since differentiation under the integral sign multiplies by $q(x)$,
\begin{equation}\label{eq:annihilated-integral}
  0
  =LU_0(z)
  =\int_A^B H(q(x))e^{zq(x)}\,dx,
\end{equation}
where
\begin{equation}\label{eq:H-def}
  H(t)=\prod_{\lambda\in\Lambda}(t-\lambda)^N
\end{equation}
is nonzero.  Because $0\in\Lambda$ was introduced to homogenize the value
relation, $H$ includes a factor $t^N$ even when $0$ is neither an endpoint
value nor a critical value; no geometric assertion about $0$ is intended.
Expanding at $z=0$ gives
\begin{equation}\label{eq:phase-moments}
  \int_A^B q(x)^mH(q(x))\,dx=0
  \qquad(m\ge0).
\end{equation}

We now rule out \eqref{eq:phase-moments} for every nonconstant polynomial phase.

\section{Phase-polynomial moment rigidity}\label{sec:exp-phase-rigidity}

\begin{lemma}[Constellation representation]\label{lem:constellation}
Let $q\in\C[x]$ be nonconstant and let $A\ne B$.  There exists a finite embedded tree $T$ in the $t$-plane, containing all finite branch values of $q$ and the endpoint values $q(A),q(B)$, together with an oriented path $\Gamma$ from $A$ to $B$ in the finite graph $q^{-1}(T)$, such that on every open edge $e$ of $T$ the path is a signed finite union of analytic inverse branches $x_i(t)$ of $q$.
\end{lemma}

\begin{proof}
Choose a regular center and pairwise disjoint analytic arcs from it to the distinct finite branch values and endpoint values, omitting repeated arms.  Their union is an embedded star, hence a finite tree.  Over the center, its preimage consists of the $n$ sheets of the covering.  The lifts of the branch arms join these sheets according to the branch-monodromy generators.  Since the polynomial covering is connected, its monodromy action is transitive, so the resulting preimage graph is connected.  Adding the endpoint arms preserves connectedness.  The graph contains $A$ and $B$, and therefore contains a path $\Gamma$ joining them.  On the interior of every edge the covering is unramified, so $\Gamma$ decomposes there into signed analytic inverse branches.  This is the standard polynomial constellation construction; see \cite[\S2.1]{PakovichMuzychuk}.
\end{proof}

\begin{lemma}[Phase-polynomial moment rigidity]\label{lem:phase-moment}
Let $q\in\C[x]$ be nonconstant, let $A\ne B$, and let $0\ne H\in\C[t]$.  Then the identities
\[
  \int_A^B q(x)^mH(q(x))\,dx=0
  \qquad(m\ge0)
\]
cannot all hold.
\end{lemma}

\begin{proof}
Let $T$ and $\Gamma$ be supplied by \cref{lem:constellation}.  In particular,
$q(\Gamma)\subseteq T$.  Subdivide the finite graph $q^{-1}(T)$ at $A$ and
$B$, and regard $\Gamma$ as an oriented finite edge path from $A$ to $B$.
For every $m\ge0$, the differential
\[
  q(x)^mH(q(x))\,dx
\]
is a polynomial differential and hence has a polynomial primitive.  Its
integral depends only on the endpoints.  Therefore the assumed identities
remain valid on $\Gamma$:
\begin{equation}\label{eq:phase-moments-Gamma}
  \int_\Gamma q(x)^mH(q(x))\,dx=0
  \qquad(m\ge0).
\end{equation}

Define, for $s\in\C\setminus T$,
\begin{equation}\label{eq:phase-Cauchy-transform}
  \mathcal C_H(s)
  =\int_\Gamma\frac{H(q(x))}{s-q(x)}\,dx.
\end{equation}
The path $\Gamma$ is compact and piecewise analytic, and the denominator in
\eqref{eq:phase-Cauchy-transform} is bounded away from zero locally uniformly
for $s\in\C\setminus T$.  Thus $\mathcal C_H$ is holomorphic on
$\C\setminus T$, with differentiation under the integral sign justified on
compact subsets.  Put
\[
  M=\max_{x\in\Gamma}|q(x)|.
\]
For $|s|>M$, the geometric series
\[
  \frac1{s-q(x)}=\sum_{m\ge0}\frac{q(x)^m}{s^{m+1}}
\]
converges uniformly on $\Gamma$.  Using \eqref{eq:phase-moments-Gamma}, we
obtain
\[
  \mathcal C_H(s)
  =\sum_{m\ge0}\frac1{s^{m+1}}
    \int_\Gamma q(x)^mH(q(x))\,dx
  =0
  \qquad(|s|>M).
\]
A finite embedded tree has connected complement in the plane.  Hence the
identity theorem gives
\begin{equation}\label{eq:phase-Cauchy-zero}
  \mathcal C_H\equiv0
  \qquad\text{on }\C\setminus T.
\end{equation}

We next express this Cauchy transform edge by edge.  Orient every open edge
$e$ of $T$.  Since all finite branch values of $q$ are vertices of $T$, the
map $q$ is unramified over $e$.  If $n=\deg q$, there are analytic inverse
branches
\[
  x_{e,1},\ldots,x_{e,n}:e\longrightarrow\C,
  \qquad q(x_{e,i}(t))=t.
\]
Let $f_{e,i}\in\Z$ be the signed number of times that $\Gamma$ traverses the
lift $x_{e,i}(e)$, relative to the chosen orientation of $e$, and set
\begin{equation}\label{eq:phase-edge-density-referee}
  \psi_e(t)
  =\sum_{i=1}^n f_{e,i}x_{e,i}'(t)
  =\sum_{i=1}^n\frac{f_{e,i}}{q'(x_{e,i}(t))}.
\end{equation}
The edge integrals below are convergent even at branch-value vertices.  Indeed,
if $c$ is an endpoint of $e$, $\xi\in q^{-1}(c)$ has ramification index
$r\ge1$, and a lifted branch tends to $\xi$, then locally
\[
  q(x)-c=(x-\xi)^r u(x),
  \qquad u(\xi)\ne0.
\]
After choosing the corresponding branch of $(t-c)^{1/r}$ along $e$, one has
\[
  x_{e,i}(t)-\xi=O\bigl((t-c)^{1/r}\bigr),
  \qquad
  x_{e,i}'(t)=O\bigl(|t-c|^{1/r-1}\bigr).
\]
Since $1/r-1>-1$, this singularity is integrable.  Changing variables
$t=q(x)$ separately on every lifted edge therefore gives the absolutely
convergent representation
\begin{equation}\label{eq:phase-edge-Cauchy-referee}
  \mathcal C_H(s)
  =\sum_e\int_e\frac{H(t)\psi_e(t)}{s-t}\,dt,
  \qquad s\in\C\setminus T.
\end{equation}

Fix an open edge $e$ and an interior point $t_0\in e$ with $H(t_0)\ne0$.
Choose a compact analytic subarc $\eta\Subset e$ containing $t_0$ in its
interior and a disk $D$ about $t_0$ such that
\[
  \overline D\cap(T\setminus\eta)=\varnothing.
\]
After shrinking $\eta$ if necessary, the inverse branches occurring in
\eqref{eq:phase-edge-density-referee} extend holomorphically to a neighborhood
of $\eta$.  Thus
\[
  \rho(t):=H(t)\psi_e(t)
\]
extends holomorphically there.  By the integrability just proved and the
positive distance from $\overline D$ to $T\setminus\eta$, the function
\[
  R(s):=\mathcal C_H(s)-\int_\eta\frac{\rho(t)}{s-t}\,dt
\]
extends holomorphically throughout $D$.  Hence the entire lateral jump of
$\mathcal C_H$ across $\eta$ is the lateral jump of the local Cauchy integral
\[
  J(s)=\int_\eta\frac{\rho(t)}{s-t}\,dt.
\]

We verify that jump directly.  Let $a,b$ be the initial and terminal endpoints
of the oriented arc $\eta$.  For $s$ near $t_0$, define
\[
  K(t,s)=
  \begin{cases}
    \dfrac{\rho(t)-\rho(s)}{t-s},&t\ne s,\\[2mm]
    \rho'(s),&t=s.
  \end{cases}
\]
The function $K$ is jointly holomorphic near $\eta\times\eta$, and
\begin{equation}\label{eq:local-Cauchy-splitting}
  J(s)
  =\rho(s)\int_\eta\frac{dt}{s-t}
   -\int_\eta K(t,s)\,dt.
\end{equation}
The second term in \eqref{eq:local-Cauchy-splitting} is holomorphic across
$\eta$.  On either side of the arc, the first scalar integral is a branch of
\[
  \int_\eta\frac{dt}{s-t}
  =\log(s-a)-\log(s-b).
\]
Its two lateral determinations differ by
$2\pi i\varepsilon_e$, where $\varepsilon_e\in\{\pm1\}$ depends only on the
orientation of $\eta$ and on which side is denoted by $+$.  Taking lateral
limits in \eqref{eq:local-Cauchy-splitting} therefore gives the local
Sokhotski--Plemelj formula
\begin{equation}\label{eq:phase-local-jump-referee}
  (\mathcal C_H)_+(t_0)-(\mathcal C_H)_-(t_0)
  =2\pi i\varepsilon_e\,H(t_0)\psi_e(t_0).
\end{equation}

Both sides of the local arc lie in $\C\setminus T$, and
\eqref{eq:phase-Cauchy-zero} says that $\mathcal C_H$ vanishes identically on
each side.  The left side of \eqref{eq:phase-local-jump-referee} is therefore
zero.  Since $H(t_0)\ne0$, we conclude that
\[
  \psi_e(t_0)=0.
\]
The point $t_0$ was arbitrary outside the finite zero set of the nonzero
polynomial $H$.  The function $\psi_e$ is analytic on the open edge and hence
continuous there, so it follows that
\begin{equation}\label{eq:phase-density-zero}
  \psi_e\equiv0
  \qquad\text{on every open edge }e\text{ of }T.
\end{equation}

Finally, decompose the line integral of $dx$ over the same lifted edges.  The
endpoint-integrability estimate above justifies the resulting improper edge
integrals, and \eqref{eq:phase-density-zero} gives
\[
  B-A
  =\int_\Gamma dx
  =\sum_e\sum_{i=1}^n f_{e,i}\int_e x_{e,i}'(t)\,dt
  =\sum_e\int_e\psi_e(t)\,dt
  =0.
\]
This contradicts $A\ne B$.  Therefore the phase-moment identities cannot all
hold.
\end{proof}

\section{Proof of the exponential-integral theorem}\label{sec:exp-main-proof}

\begin{proof}[Proof of \cref{thm:exp-integral}]
Suppose first that $q(x)=ax+b$ with $a\ne0$, and put
\[
  \alpha=aB+b,
  \qquad
  \beta=aA+b.
\]
Then $\alpha\ne\beta$ are algebraic and
\[
  \int_A^B e^{q(x)}\,dx
  =\frac{e^\alpha-e^\beta}{a}.
\]
If this value were an algebraic number $\gamma$, then
\[
  e^\alpha-e^\beta-a\gamma e^0=0.
\]
After combining terms if one of $\alpha,\beta$ equals $0$, this is a nontrivial algebraic linear relation among exponentials of distinct algebraic numbers drawn from $\{\alpha,\beta,0\}$.  The Lindemann--Weierstrass theorem excludes such a relation; see \cite[Chapter~1]{Baker}.  Hence the integral is transcendental.

Now let $n=\deg q\ge2$.  Apply the algebraic affine normalization of \cref{eq:normalized-q}; it suffices to treat the normalized polynomial, and we retain the notation $q,A,B$.  Let
\[
  U_0(z)=\int_A^B e^{zq(x)}\,dx.
\]
Assume, toward a contradiction, that $U_0(1)$ is algebraic.  By \cref{prop:value-to-exp-poly}, $U_0$ has the exponential-polynomial representation \eqref{eq:exp-polynomial-U0}.  Sections \ref{sec:exp-annihilation} and \ref{sec:exp-phase-rigidity} then give a nonzero polynomial $H$ satisfying the impossible moment identities \eqref{eq:phase-moments}.  This contradiction proves that $U_0(1)$ is transcendental.
\end{proof}



\section{Completion of the proof of the Factorial Conjecture}\label{sec:FC-proof}

\begin{proof}[Proof of \cref{thm:FC2}]
Assume that a nonzero counterexample exists.  By
\cref{prop:algebraic-specialization}, we may take
$f\in\Qbar[x,y]$.  A nonzero constant has nonzero first factorial moment,
so $D=\deg f\ge1$.

If $A_0(t)=f_D(t,1-t)$ is nonconstant, then
\cref{thm:nonflat-exclusion} gives infinitely many nonzero factorial
moments, a contradiction.  Hence $A_0$ is constant, and
\cref{cor:flat-top-reduction} gives
\[
  f_D=a(x+y)^D,
  \qquad a\in\Qbar^\times.
\]
Multiplying $f$ by $a^{-1}$ preserves the vanishing of all positive power
moments, so we may assume $f_D=(x+y)^D$.

Put
\[
  q(t)=\frac{A_1(t)}D
  =\frac{f_{D-1}(t,1-t)}D\in\Qbar[t].
\]
By \cref{prop:first-flat-symbol},
\begin{equation}\label{eq:final-zero-integral}
  \int_0^1 e^{q(t)}\,dt=0.
\end{equation}
If $q$ is constant, the left side is the nonzero number $e^q$.  If $q$ is
nonconstant, \cref{thm:exp-integral} with $A=0$ and $B=1$ says that the
left side is transcendental and therefore nonzero.  Both cases contradict
\eqref{eq:final-zero-integral}.  Thus no nonzero counterexample exists.
\end{proof}

\section*{Acknowledgments and AI-assisted research disclosure}

This manuscript was developed through interactive work between Christopher
D. Long and ChatGPT 5.6 Sol.  ChatGPT 5.6 Sol assisted in proof discovery,
organization, symbolic checking, reference verification, adversarial
auditing, and drafting.  Claude Opus 5 contributed the semisimple-projector
strategy, the reduction to phase-polynomial moments, and extensive
adversarial audits.  The AI systems are not authors.  The human author bears
full responsibility for all statements, proofs, citations, and any remaining
errors.  This draft has not yet undergone independent peer review or formal
verification.

\footnotesize
\begin{thebibliography}{99}

\bibitem{Baker}
A.~Baker,
\emph{Transcendental Number Theory},
Cambridge Mathematical Library,
Cambridge University Press, Cambridge, 1990.

\bibitem{Beukers}
F.~Beukers,
\emph{A refined version of the Siegel--Shidlovskii theorem},
Ann. of Math. (2) \textbf{163} (2006), no.~1, 369--379.
\href{https://doi.org/10.4007/annals.2006.163.369}{doi:10.4007/annals.2006.163.369}.

\bibitem{deCataldoMigliorini}
M.~A.~A. de~Cataldo and L.~Migliorini,
\emph{The decomposition theorem, perverse sheaves and the topology of algebraic maps},
Bull. Amer. Math. Soc. (N.S.) \textbf{46} (2009), no.~4, 535--633.
\href{https://doi.org/10.1090/S0273-0979-09-01260-9}{doi:10.1090/S0273-0979-09-01260-9}.


\bibitem{HTT}
R.~Hotta, K.~Takeuchi, and T.~Tanisaki,
\emph{$D$-Modules, Perverse Sheaves, and Representation Theory},
Progress in Mathematics, vol.~236,
Birkh\"auser Boston, Boston, MA, 2008.
\href{https://doi.org/10.1007/978-0-8176-4523-6}{doi:10.1007/978-0-8176-4523-6}.

\bibitem{LiuSun}
D.~Liu and X.~Sun,
\emph{The Factorial Conjecture and images of locally nilpotent derivations},
Bull. Aust. Math. Soc. \textbf{101} (2020), no.~1, 71--79.
\href{https://doi.org/10.1017/S0004972719000546}{doi:10.1017/S0004972719000546}.

\bibitem{PakovichMuzychuk}
F.~Pakovich and M.~Muzychuk,
\emph{Solution of the polynomial moment problem},
Proc. Lond. Math. Soc. (3) \textbf{99} (2009), no.~3, 633--657.
\href{https://doi.org/10.1112/plms/pdp010}{doi:10.1112/plms/pdp010}.


\bibitem{vdEWZ}
A.~van den Essen, D.~Wright, and W.~Zhao,
\emph{On the Image Conjecture},
J. Algebra \textbf{340} (2011), 211--224.
\href{https://doi.org/10.1016/j.jalgebra.2011.04.036}{doi:10.1016/j.jalgebra.2011.04.036}.

\end{thebibliography}

\end{document}
